\ExerciseGroupsStart
\ExerciseSection

\begin{enumerate}
\def\labelenumi{\arabic{enumi})}
\tightlist
\item
  一个正实值函数 \(f\) （定义在 \(\mathbf{X} \times \mathbf{X}\) 上）是 \(\mathbf{X}\) 上的伪度量，当且仅当 \(f(x, x) = 0\) 对所有 \(x \in \mathbf{X}\) 成立，并且
\end{enumerate}

\[
f (x, y) \leqslant f (x, z) + f (y, z)
\]

对所有 \(x, y, z\) ∈ \(\mathbf{X}\) 成立 。

\begin{enumerate}
\def\labelenumi{\arabic{enumi})}
\setcounter{enumi}{1}
\tightlist
\item
  设 \(f\) 是 \(\mathbf{X} \times \mathbf{X}\) 到 \([0, +\infty]\) 中的一个映射。那么，由诸集 \(\overline{f}^{1}([0, a])\) （其中 \(a\) 取遍所有实数 \(>0\) ）组成的族，构成 \(\mathbf{X}\) 上某个一致结构的围域基本系，当且仅当 \(f\) 满足下列条件：\(a) f(x, x) = 0\) 对所有 \(x \in \mathbf{X}\) 成立；\(b)\) 对每个 \(a > 0\) ，存在 \(b > 0\) ，使得关系 \(f(x, y) \leqslant b\) 蕴含 \(f(y, x) \leqslant a\) ；\(c)\) 对每个 \(a > 0\) ，存在 \(c > 0\) ，使得关系 \(f(x, z) \leqslant c\) 与 \(f(z, y) \leqslant c\) 蕴含 \(f(x, y) \leqslant a\) 。
\end{enumerate}

条件 \(b)\) 与 \(c)\) 特别在下述情形得到满足：存在 \(\varphi\) ——一个 \(I = [0, +\infty]\) 到其自身上的连续映射，它在点 \(0\) 处取值零——以及 \(\psi\) ——一个 \(I \times I\) 到 \(I\) 中的连续映射，它在点 \((0, 0)\) 处取值零——使得恒有

\[
f (y, x) \leqslant \varphi (f (x, y)) \quad \text { and } \quad f (x, y) \leqslant \psi (f (x, z), f (z, y)).
\]

\begin{enumerate}
\def\labelenumi{\arabic{enumi})}
\setcounter{enumi}{2}
\tightlist
\item
  设 X 是一个拓扑空间，C 是 X 的拓扑，\((f_{t})\) 是 X 上伪度量的一个饱和族（第 2 号），U 是由族 \((f_{t})\) 定义的一致结构。
\end{enumerate}

\begin{enumerate}
\def\labelenumi{\alph{enumi})}
\item
  由 \(\mathcal{U}\) 诱导的拓扑粗于 \(\mathfrak{C}\) ，当且仅当诸 \(f_{t}\) 在 \(\mathbf{X} \times \mathbf{X}\) 上连续（关于拓扑 \(\mathfrak{C}\) 与其自身的乘积）。
\item
  由 \(\mathcal{U}\) 诱导的拓扑细于 \(\mathcal{C}\) ，当且仅当对每个 \(x_0 \in \mathbf{X}\) 与每个邻域 \(\mathbf{V}\) ——它包围 \(x_0\) （按拓扑 \(\mathcal{C}\) ）——，存在一个指标 \(\iota\) 与一个实数 \(a > 0\) ，使得 \(f_{\iota}(x_0, x) \geqslant a\) 对所有 \(x \in \mathbb{C}\mathbf{V}\) 成立 。
\end{enumerate}

\begin{enumerate}
\def\labelenumi{\arabic{enumi})}
\setcounter{enumi}{3}
\tightlist
\item
  设 X 是一个非豪斯多夫的可一致化空间，R 是关系`` \(y \in \overline{\{x\}}\) ''（在 X 的两点 \(x, y\) 之间）。
\end{enumerate}

\begin{enumerate}
\def\labelenumi{\alph{enumi})}
\item
  证明 R 是 X 上的一个等价关系，X 到豪斯多夫空间中的每个连续映射都与关系 R 相容（《集合论》，R，§ 5，第 7 号），并且商空间 X/R 是完全正则的。
\item
  若 \(\mathcal{U}\) 是任一与 X 的拓扑相容的一致结构，则与 \(\mathcal{U}\) 相伴的豪斯多夫一致结构（第 II 章，§ 3，第 9 号）定义在商空间 X/R 上，且与 X/R 的拓扑相容。拓扑空间 X/R 称为与可一致化空间 X 相伴的完全正则空间。
\end{enumerate}

\begin{enumerate}
\def\labelenumi{\arabic{enumi})}
\setcounter{enumi}{4}
\tightlist
\item
  设 X 是一个可一致化空间。证明 X 上所有在 \(X \times X\) 上连续的伪度量构成的族定义了 X 上一个与 X 的拓扑相容的一致结构。这个一致结构 \(U_{0}\) 称为 X 上的万有一致结构；它是与 X 的拓扑相容的所有一致结构中最细的一个。若 Y 是任一一致空间，且 \(f: X \to Y\) 是连续映射，则 f 关于 X 上的万有一致结构是一致连续的；对于 X 上与其拓扑相容的任何其他一致结构，情形并非如此。若存在 X 上的一致结构 U ，它与 X 的拓扑相容，且使 X 赋有该一致结构后为完备空间，则 X 关于任何一致结构 \(U'\) （它比 U 更细且比 \(U_{0}\) 更粗）也是完备空间。（注意，柯西滤子关于 \(U'\) 与关于 \(U_{0}\) 是同样的。）
\end{enumerate}

¶ 6) a) 设 X 是一个完全正则空间，U 是一个与 X 的拓扑相容的一致结构。设 U* 是 X 上使所有关于 U 一致连续的 X 到 {[}0, 1{]} 中的映射都一致连续的最粗一致结构。证明一致结构 U* 是豪斯多夫的且与 X 的拓扑相容，并且 X 关于 U* 是预紧的。

\begin{enumerate}
\def\labelenumi{\alph{enumi})}
\setcounter{enumi}{1}
\tightlist
\item
  设 K 是一个紧空间。证明每个关于 U 一致连续的映射 \(f: X \to K\) 关于 \(U^{*}\) 也一致连续（注意，K 上唯一的一致结构是使 K 到区间 {[}0, 1{]} 中的所有连续映射都一致连续的最粗一致结构）。由此证明 \(U^{*}\) 是 X 上比 U 更粗且使 X 关于它为预紧的一致结构中最细的一个。
\end{enumerate}

\begin{enumerate}
\def\labelenumi{\arabic{enumi})}
\setcounter{enumi}{6}
\tightlist
\item
  设 X 是一个完全正则空间。若 \(U_{0}\) 表示 X 上的万有一致结构（习题 5），则 \(U_{0}^{*}\) （习题 6）是 X 上使所有 X 到 {[}0, 1{]} 中的连续映射都一致连续的最粗一致结构。设 \(\beta X\) 表示关于一致结构 \(U_{0}^{*}\) 完备化 X 所得的紧空间。\(\beta X\) 称为 X 的斯通–切赫紧化。
\end{enumerate}

\begin{enumerate}
\def\labelenumi{\alph{enumi})}
\item
  证明 X 到紧空间 K 中的每个连续映射都可以扩张为 \(\beta X\) 到 K 中的连续映射（参看《集合论》，第 IV 章，§ 3）。
\item
  设 \(f\) 是 \(\mathbf{X}\) 到一个稠密子集 \(\mathbf{X}'\) （紧空间 \(\mathbf{K}\) 的）上的同胚，\(\overline{f}\) 是连续映射 \(\beta \mathbf{X} \to \mathbf{K}\) （它扩张 \(f\) ）。证明 \(\overline{f}(\beta \mathbf{X} - \mathbf{X}) = \mathbf{K} - \mathbf{X}'\) 。{[}假设存在一点 \(a \in \beta \mathbf{X} - \mathbf{X}\) 使 \(\overline{f}(a) = f(b)\) ，其中 \(b \in \mathbf{X}\) ；考察连续映射 \(g: \beta \mathbf{X} \to [0, 1]\) ，它满足 \(g(a) = 1\) 与 \(g(b) = 0\) ，并令 \(P\) 表示所有 \(x \in \beta X\) （满足 \(g(x) \geqslant 1/2\) ）所成的集。存在连续映射 \(h: K \to [0, 1]\) ，使 \(h(f(b)) = 0\) 且 \(h(f(x)) = 1\) 对所有 \(x \in P\) 成立；证明关系 \(h(f(a)) = 0\) 与 \(a\) 属于 \(P\) 在 \(\beta X\) 中的闭包这一事实相矛盾。{]}
\end{enumerate}

\begin{enumerate}
\def\labelenumi{\arabic{enumi})}
\setcounter{enumi}{7}
\tightlist
\item
  设 X 是一个完全正则空间，\(\beta X\) 是它的斯通–切赫紧化。
\end{enumerate}

\begin{enumerate}
\def\labelenumi{\alph{enumi})}
\item
  X 上的滤子 \(\mathfrak{F}\) 称为完全正则滤子，如果 \(\mathfrak{F}\) 有一个由开集组成的基 \(\mathfrak{B}\) ，使得对每个集 \(A \in \mathfrak{B}\) ，存在含于 A 的集 \(B \in \mathfrak{B}\) ，以及一个连续映射 \(f: X \to [0, 1]\) ，它在 B 上等于 0 、在 \(\mathbb{C}A\) 上等于 1 。完全正则滤子称为极大的，如果不存在严格更细的完全正则滤子。证明：若 \(\mathfrak{F}\) 是 E 上任一完全正则滤子，则存在一个比 \(\mathfrak{F}\) 更细的极大完全正则滤子（用佐恩引理）。
\item
  完全正则滤子 \(\mathfrak{F}\) 是极大的，当且仅当：对 X 的每对开集 A, B ，其中 \(B \subset A\) 且存在 X 到 \([0, 1]\) 中的连续映射在 B 上等于 0 、在 \(\complement A\) 上等于 1 ，则或者 \(A \in \mathfrak{F}\) ，或者 \(A \notin \mathfrak{F}\) 而 \(\mathfrak{F}\) 中有一个不与 B 相交的集（若 \(\mathfrak{F}\) 的所有集都与 B 相交，考察由 \(\mathfrak{F}\) 的诸集与诸集 \(\overline{f}([0, a[)\) （其中 \(a \in ]0, 1[\) ）生成的滤子）。
\item
  证明每个极大完全正则滤子 \(\mathfrak{F}\) 都是 \(\mathbf{X}\) 上关于由 \(\beta\mathbf{X}\) 的一致结构诱导的一致结构的柯西滤子{[}证明每个连续映射 \(f: \mathbf{X} \to [\mathrm{0}, \mathrm{1}]\) 关于 \(\mathfrak{F}\) 有极限（用反证法，并利用 \(b)\) ）{]}。
\end{enumerate}

\(d)\) 两个不同的极大完全正则滤子 \(\mathfrak{F},\mathfrak{F}^{\prime}\) 不能收敛到 \(\beta\mathbf{X}\) 的同一点{[}注意，由 \(b)\) ，存在开集 \(\mathbf{A}\in\mathfrak{F}\) 与开集 \(\mathbf{A}^{\prime}\in\mathfrak{F}^{\prime}\) ，使 \(\mathbf{A}\cap\mathbf{A}^{\prime}=\varnothing\) ；然后利用公理 \((\mathrm{O}_{\mathbf{IV}})\) 进行反证{]}。

\begin{enumerate}
\def\labelenumi{\alph{enumi})}
\setcounter{enumi}{4}
\tightlist
\item
  任一点 \(x_{0} \in \beta X\) 的邻域滤子在 X 上的迹是一个极大完全正则滤子 B （用反证法，并利用完全正则滤子的定义以及 \(\beta X\) 中一点的邻域滤子的定义）。
\end{enumerate}

\begin{enumerate}
\def\labelenumi{\arabic{enumi})}
\setcounter{enumi}{8}
\tightlist
\item
  \begin{enumerate}
  \def\labelenumii{\alph{enumii})}
  \tightlist
  \item
    设 X 是一个拓扑空间，\(\mathcal{G}\) 是它的拓扑，\((f_{\iota})_{\iota \in \mathbb{I}}\) 是 X 到 K = {[}0, 1{]} 中的一族映射，它们在拓扑 \(\mathcal{G}\) 中连续。设 \(\mathcal{G}_0\) 是 X 上使所有 \(f_{\iota}\) 都连续的最粗拓扑。证明 \(\mathcal{G}_0 = \mathcal{G}\) ，只要诸集 \(\overline{f}_{\iota}^{1}([0, a[)\) （其中 \(\iota \in I\) 且 \(a \in ]0, 1[,\) ）组成的族是 \(\mathcal{G}\) 的一个子基（第 I 章，§ 2，第 3 号）。此时空间 X 是可一致化的，并且若它是豪斯多夫的，它就同胚于方体 K \(^{\mathbf{I}}\) 的一个子空间。
  \end{enumerate}
\end{enumerate}

\begin{enumerate}
\def\labelenumi{\alph{enumi})}
\setcounter{enumi}{1}
\tightlist
\item
  假设族 \((f_{t})\) 是 X 到 K 中的所有连续映射（关于 C）构成的族。证明：若 Y 是任一紧空间，则每个关于 C 连续的 X 到 Y 中的映射关于 \(C_{0}\) 也连续（把 Y 嵌入一个方体）。拓扑 C 是可一致化的，当且仅当 \(C_{0} = C\) 。
\end{enumerate}

\begin{enumerate}
\def\labelenumi{\arabic{enumi})}
\setcounter{enumi}{9}
\tightlist
\item
  设 X 是一个完全正则空间，K 是 X 的一个紧子集，V 是 K 在 X 中的一个邻域。
\end{enumerate}

\begin{enumerate}
\def\labelenumi{\alph{enumi})}
\item
  证明存在 X 到 {[}0, 1{]} 中的连续映射，它在 K 上等于 1 ，在 \(\complement V\) 上等于 0 {[}利用 \((\mathrm{O}_{\mathbf{IV}})\) ，用有限多个适当选取的邻域覆盖 K{]}。
\item
  设 \(X'\) 是把 K 的所有点等同起来所得到的 X 的商空间。证明 \(X'\) 是完全正则的。
\end{enumerate}

¶ 11) 设 X 是一个完全正则空间。X 中两个闭集 A, B 称为完全分离的，如果存在 X 到 {[}0, 1{]} 中的连续映射 f ，使得在 A 上 \(f(x) = 0\) ，在 B 上 \(f(x) = 1\) 。a) 设 \(\beta X\) 是 X 的斯通–切赫紧化。证明 X 中两个闭集 A 与 B 完全分离，当且仅当它们在 \(\beta X\) 中的闭包不相交{[}利用习题 7 a) 与 10 a){]}。

\begin{enumerate}
\def\labelenumi{\alph{enumi})}
\setcounter{enumi}{1}
\tightlist
\item
  设 \(h\) 是 \(\mathbf{X}\) 到一个稠密子集 \(\mathbf{X}'\) （紧空间 \(\mathbf{K}\) 的）上的同胚，\(\overline{h}\) 是连续映射 \(\beta \mathbf{X} \to \mathbf{K}\) （它扩张 \(h\) ）。假设对每对完全分离的闭子集 \(\mathbf{A}, \mathbf{B}\) （属于 \(\mathbf{X}\) ），\(\overline{h}(\mathbf{A})\) 与 \(h(\mathbf{B})\) 在 \(\mathbf{K}\) 中的闭包不相交。证明 \(\overline{h}\) 是 \(\beta \mathbf{X}\) 到 \(\mathbf{K}\) 上的一个同胚。（证明 \(\overline{h}\) 是单射：若 \(a, b\) 是 \(\beta \mathbf{X}\) 中满足 \(\overline{h}(a) = \overline{h}(b)\) 的两个不同点，考察连续映射 \(f: \beta \mathbf{X} \to [\mathbf{0}, \mathbf{1}]\) ，它满足 \(f(a) = 0\) 、\(f(b) = 1\) ，并考察集 \(X \cap \overline{f}[0, 1/3]\) 与 \(X \cap \overline{f}[2/3, 1]\) 。）
\end{enumerate}

¶ 12) 设 X 是一个完全正则空间，βX 是它的斯通–切赫紧化。

\begin{enumerate}
\def\labelenumi{\alph{enumi})}
\item
  设 \(f\) 是 \(\beta X\) 上的一个有限连续实值函数，使 \(f(x) > 0\) 对所有 \(x \in X\) 成立，且使集 \(A = \overline{f}^1(0) \subset \beta X - X\) 非空。那么存在点的序列 \((a_n)\) （属于 \(X\) ），使数列 \(\lambda_n = f(a_n) > 0\) 严格递减，且 \(\lim_{n \to \infty} \lambda_n = 0\) 。对每个整数 \(n > 0\) ，设 \(I_n\) 是以 \(\lambda_n\) 为中心的开区间（在 \(R_+^*\) 中），使得闭区间 \(\overline{I}_n\) 两两不相交，并令 \(M_n = \overline{f}^1(I_n) \cap X\) 。对每个子集 \(H\) ⊂ \(N\) ，令 \(M_H\) 为诸 \(M_n\) （其中 \(n \in H\) ）的并；对每个滤子 \(\mathfrak{F}\) （在 \(N\) 上且比弗雷歇滤子更细），令 \(\mathfrak{F}'\) 表示滤子（在 \(X\) 上），它以诸 \(M_H(H \in \mathfrak{F})\) 为基。证明 \(\mathfrak{F}'\) 是完全正则滤子（习题 8），并且若 \(\mathfrak{F}_1, \mathfrak{F}_2\) 是两个不同的超滤子（在 \(N\) 上），其中每个都比弗雷歇滤子更细，则 \(X\) 上不存在比 \(\mathfrak{F}_1', \mathfrak{F}_2'\) 中每个都更细的滤子。由此推出 Card \((A) \geqslant 2^{2\text{Card}(N)}\) {[}利用习题 8 d) 与第 I 章，§ 4，习题 5{]}。
\item
  假设 X 是无限且离散的。证明由 \(\beta\mathbf{X}\) 的一致结构在 X 上诱导的一致结构是有限分划的一致结构（第 II 章，§ 2，第 2 号及 § 4，习题 12）{[}利用习题 11 b){]}。由此证明 Card \((\beta\mathbf{X}) = 2^{2^{\operatorname{Card}(\mathbb{N})}}\) （参看第 I 章，§ 4，习题 5）。证明 X 上存在连续实值函数 \(f \geqslant 0\) ，使 \(\overline{f}^{1}(0)\) 非空且含于 \(\beta\mathbf{X} - \mathbf{X}\) 。
\item
  在与 b) 相同的假设下，证明：若 A 是 \(\beta X\) 的一个无限闭子集，则 \(\operatorname{Card}(A) \geqslant 2^{2\operatorname{Card}(N)}\) 。{[}首先证明存在点的无限序列 \((a_{n})\) （属于 A ），且对每个 n 存在邻域 \(V_{n}\) （\(a_{n}\) 在 \(\beta X\) 中的），诸 \(V_{n}\) 两两不相交；由此推出定义在诸 \(a_{n}\) 所成的集 D 上的每个有界实值函数 f 都可以扩张为 \(\beta X\) 上的连续函数；为此，考察定义在 X 上的实值函数 g ，使 \(f(a_{n})\) 是它在 \(X \cap V_{n}\) 每点处的值。于是断言 \(\overline{D} \subset A\) 同胚于 D 的斯通–切赫紧化（参看 § 4，习题 17）。{]}
\end{enumerate}

\begin{enumerate}
\def\labelenumi{\arabic{enumi})}
\setcounter{enumi}{12}
\item
  设 G 是一个局部紧但非紧的拓扑群。证明：在乘积空间 \(H = G \times G\) 上，由其斯通–切赫紧化 \(\beta H\) 的一致结构诱导的一致结构，严格细于由 \((\beta G) \times (\beta G)\) 的一致结构诱导的一致结构。{[}假设结论不真，把习题 7 a) 应用于 H 到 G 中的映射 \((x, y) \to xy^{-1}\) ，证明 G 将同构于一个紧群的子群；然后应用第 III 章，§ 3，第 3 号，命题 4 的推论 1 。{]}
\item
  若拓扑空间 X 的每一点都有一个闭邻域，它是 X 的可一致化子空间，证明 X 是可一致化的。
\end{enumerate}

¶ 15) a) 设 X 是一个局部紧空间，\(\Phi\) 是所有满足下列条件的连续映射 \(g: \mathbf{X} \to [\mathfrak{0}, \mathfrak{1}]\) 组成的集：g 在某个紧集（依赖于 \(g\) ）的余集上为零。设 \(\mathcal{U}_1\) 是 X 上使所有映射 \(g \in \Phi\) 都一致连续的最粗一致结构。证明 \(\mathcal{U}_1\) 与 X 的拓扑相容，X 的完备化 \(\hat{\mathbf{X}}\) （关于 \(\mathcal{U}_1\) ）是紧的，且 \(\mathbf{X} - \hat{\mathbf{X}}\) 至多由一点组成（证明：若 X 不是紧的，则 X 的相对紧子集的余集所成的滤子是关于 \(\mathcal{U}_1\) 的柯西滤子）。

\begin{enumerate}
\def\labelenumi{\alph{enumi})}
\setcounter{enumi}{1}
\item
  证明 \(\mathfrak{U}_1\) 是与 \(\mathbf{X}\) 的拓扑相容的最粗一致结构。
\item
  反之，设 X 是一个完全正则空间，使得与 X 的拓扑相容的一致结构组成的集有一个最粗元 \(U_{1}\) 。证明 X 是局部紧的。（利用习题 6 证明 X 关于 \(U_{1}\) 是预紧的；然后证明 X 在其关于 \(U_{1}\) 的完备化中的余集不可能多于一点。）
\item
  假设 X 是局部紧且 \(\sigma\) 紧的，因而是递增序列 \((\mathbf{U}_n)\) （由相对紧开集组成，且满足 \(\overline{\mathbf{U}}_n \subset \mathbf{U}_{n+1}\) ）之并（第 I 章，§ 9，第 9 号，命题 15）。证明存在一个连续实值函数 \(f\) （定义在 \(\mathbf{X}\) 上），使得 \(f(x) \leqslant n\) 对 \(x \in \overline{\mathbf{U}}_n\) 成立，且 \(f(x) \geqslant n\) 对 \(x \in \mathbb{C} \overline{\mathbf{U}}_n\) 成立{[}参看习题 10 a){]}。设 \(\mathcal{U}_2\) 是 \(\mathbf{X}\) 上使 \(f\) 与所有函数 \(g \in \Phi\) 都一致连续的最粗一致结构。证明 \(\mathcal{U}_2\) 与 \(\mathbf{X}\) 的拓扑相容，且存在一个围域 \(\mathbf{V}\) （属于 \(\mathcal{U}_2\) ），使得 \(\mathbf{V}(x)\) 在 \(\mathbf{X}\) 中相对紧对所有 \(x \in \mathbf{X}\) 成立（参看第 II 章，§ 4，习题 9）。
\end{enumerate}

¶ 16) 设 X 是一个完备的豪斯多夫一致空间，U 是 X 的一致结构。

\begin{enumerate}
\def\labelenumi{\alph{enumi})}
\tightlist
\item
  设 A 是 X 的一个子集，它是一个闭集序列 \((\mathbf{F}_{n})\) 之并，且 \(x \notin A\) 。证明 X 上存在连续的有限实值函数 \(f \geqslant 0\) ，使 \(f(x) = 0\) 且 \(f(y) > 0\) 对所有 \(y \in A\) 成立。（若 \(g_{n}\) 是 X 到 {[}0, 1{]} 中的连续映射，满足 \(g_{n}(x) = 0\) 且 \(g_{n}(y) = 1\) 对所有 \(y \in F_{n}\) 成立，考察函数
\end{enumerate}

\[
f (x) = \sum_ {n = 0} ^ {\infty} g _ {n} (x) / 2 ^ {n}.
\]

\begin{enumerate}
\def\labelenumi{\alph{enumi})}
\setcounter{enumi}{1}
\tightlist
\item
  设 \(h \geqslant 0\) 是 \(\mathbf{X}\) 上的连续实值函数，\(\mathbf{V}\) 是开集 \(\overline{h}^{1}(]0, +\infty[)\) 。设 \(\mathcal{U}'\) 是 \(\mathbf{V}\) 上由 U 诱导的一致结构与由伪度量
\end{enumerate}

\[
r (x, y) = \left| \frac {\mathrm{1}}{h (x)} - \frac {\mathrm{1}}{h (y)} \right|.
\]

定义的一致结构的最小上界。证明 V 关于 \(U'\) 是完备的。

\begin{enumerate}
\def\labelenumi{\alph{enumi})}
\setcounter{enumi}{2}
\tightlist
\item
  设 B 是 X 的一个子集，它是一族集 \((\mathbf{A}_{\lambda})\) 之交，其中每个集都是可数闭集族之并。证明 B 上存在一个一致结构，它与 X 的拓扑在 B 上诱导的拓扑相容，且 B 关于这个一致结构是完备的{[}利用 b){]}。
\end{enumerate}

\begin{enumerate}
\def\labelenumi{\arabic{enumi})}
\setcounter{enumi}{16}
\item
  设 X 是一个拓扑空间，它的每个点 \(x \in X\) 都有一个既开又闭的邻域组成的基本邻域系。证明 X 是可一致化的。
\item
  设 X 是一个可数的完全正则空间。证明对每个 \(x \in X\) ，x 的既开又闭的邻域构成 x 的一个基本邻域系（参看 § 6，第 4 号）。
\item
  设 X 是一个局部紧空间。证明 X 上每个下半连续的函数 \(f \geqslant 0\) 都是一族连续函数 \(\geqslant 0\) 的上包络，这族函数中的每一个都在某个紧集的余集上为零。
\item
  \begin{enumerate}
  \def\labelenumii{\alph{enumii})}
  \tightlist
  \item
    设 X 是一个完全正则空间，a 是 X 的一点。那么，存在 X 上的连续函数 \(f \geqslant 0\) ，使 \(f(a) = 0\) 且 \(f(x) > 0\) 对每个 \(x \neq a\) 成立，当且仅当存在 a 在 X 中的邻域序列 \((\mathbf{V}_n)\) ，使 \(\bigcap_{n} V_n = \{a\}\) {[}参看习题 16 a){]}。
  \end{enumerate}
\end{enumerate}

\begin{enumerate}
\def\labelenumi{\alph{enumi})}
\setcounter{enumi}{1}
\tightlist
\item
  设 X 是一个完全正则空间，它的每个点都有可数的基本邻域系。证明在斯通–切赫紧化 \(\beta\mathrm{X}\) 中，集 X 等于 \(\beta\mathrm{X}\) 中具有可数基本邻域系的点所成的集{[}利用 a) 与习题 12{]}。设 \(\mathbf{X}'\) 是另一个完全正则空间，它的每个点都有可数的基本邻域系。证明：若斯通–切赫紧化 \(\beta\mathrm{X}\) 与 \(\beta\mathrm{X}'\) 同胚，则 X 与 \(\mathbf{X}'\) 同胚。
\end{enumerate}

¶ 21) a) 设 X 是拓扑空间。证明下列两个性质等价：α) X 上每个连续有限实值函数都有界；β) X 上每个有界的连续实值函数都达到它的界。若 X 具有这两个性质，则称 X 为伪紧的。若 X 伪紧，且 f 是 X 到拓扑空间 \(X'\) 中的任一连续映射，则 \(f(X)\) 伪紧。

\begin{enumerate}
\def\labelenumi{\alph{enumi})}
\setcounter{enumi}{1}
\item
  考虑拓扑空间 X 的下列两个性质：\(\gamma\) 若 \((\mathbf{U}_n)\) 是 X 的任一可数开覆盖，则 X 是有限个闭包 \(\overline{\mathbf{U}}_n\) 的并；\(\delta\) X 上每个由开集组成的可数滤子基至少有一个触点。证明 \(\gamma\) 与 \(\delta\) 等价，且它们蕴含 X 伪紧。特别地，每个绝对闭的拓扑空间（第 I 章，§ 9，习题 19）都是伪紧的。若 X 具有性质 \(\gamma\) 与 \(\delta\) ，则 X 的每个作为 X 的开集之闭包的子空间也具有这两性质{[}参看 § 4，习题 26 b){]}。
\item
  考虑拓扑空间 \(\mathbf{X}\) 的下列性质：\(\zeta)\) \(\mathbf{X}\) 的每个局部有限开覆盖都是有限的。证明 \((\zeta) \Longrightarrow \gamma)\) 。{[}若 \((\mathbf{U}_n)\) 是 \(\mathbf{X}\) 的开子集组成的一个递增序列，诸子集覆盖 \(\mathbf{X}\) 且满足 \(\mathbf{U}_{n+1} \notin \overline{\mathbf{U}}_n\) ，考虑由诸集 \(\mathbf{U}_{n+1} \cap \mathbb{C}\overline{\mathbf{U}}_n\) 以及某个序列 \((a_n)\) 的余集所组成的覆盖，该序列满足
\end{enumerate}

\[
a _ {n + 1} \in \mathrm{U} _ {n + 1} \cap \left[ \overline {{\mathrm{U}}} _ {n}. \right]
\]

若 X 是正则的，证明 \(\gamma\Longrightarrow\zeta\) 。{[}设 \((\mathrm{U}_{\alpha})\) 是 X 的一个无限的局部有限开覆盖。用归纳法定义指标的序列 \((\alpha_{n})\) 、X 的点的序列 \((x_{n})\) ，并对每个 \(x_{n}\) 定义两个开邻域 \(V_{n}\) 与 \(W_{n}\) （\(x_{n}\) 的），使得：(i) \(\overline{V}_{n}\subset W_{n}\) ，\(W_{n}\subset U_{\alpha_{n}}\) ，且 \(W_{n}\) 只与有限个集 \(U_{\alpha}\) 相交；(ii) \(U_{\alpha_{n}}\) 不与任何 \(W_{k}\) （指标 k\textless n 者）相交。现在考虑由诸 \(W_{n}\) 与诸 \(\overline{V}_{n}\) 的并的余集所组成的覆盖。{]}

\begin{enumerate}
\def\labelenumi{\alph{enumi})}
\setcounter{enumi}{3}
\item
  在完全正则空间 X 中，性质 \(\alpha\) 、\(\beta\) 、\(\gamma\) 、\(\delta\) ) 与 \(\zeta\) 都等价，并且都与下列性质等价：\(\theta\) X 在任何与它的拓扑相容的一致结构中都是预紧的。{[}为证 \(\alpha) \Longrightarrow \zeta\) )，注意：若 \((\mathbf{U}_n)\) 是 X 的一个可数无限的局部有限开覆盖，且 \(a_n \in \mathbf{U}_n\) ，则 X 上存在连续实值函数 \(f\) ，使得 \(f(a_n) = n\) 对每个 \(n\) 成立{[}利用 \((\mathrm{O}_{\mathrm{IV}})]\) 。为证 \(\theta) \Longrightarrow \alpha\) )，利用习题 5，并注意在预紧空间中每个一致连续实值函数都有界。为证 \(\gamma) \Longrightarrow \theta\) )，注意：若 X 关于某个一致结构 \(\mathcal{U}\) 不是预紧的，则存在 \(\mathcal{U}\) 的一个对称围域 V 与 X 的点的序列 \((x_n)\) ，使得诸集 \(\mathrm{V}(x_n)\) 中任何两个都不相交。{]}
\item
  若 X 是完全正则的伪紧空间，且关于某个与它的拓扑相容的一致结构是完备的，则 X 是紧的。
\item
  在完全正则的伪紧空间 X 上，万有一致结构（习题 5）由 X 的斯通–切赫紧化的一致结构所诱导，并且是使 X 到 {[}0, 1{]} 的所有连续映射都一致连续的、唯一与 X 的拓扑相容的一致结构。
\end{enumerate}

\begin{enumerate}
\def\labelenumi{\arabic{enumi})}
\setcounter{enumi}{21}
\tightlist
\item
  设 X 是豪斯多夫一致空间，Y 是 X 的闭子集，f 是 Y 上有界的一致连续实值函数。证明 f 有一个扩张 \(\overline{f}\) 到 X 上，它在 X 上有界且一致连续。{[}我们可以假设 \(f(Y) \subset [0, 1]\) 。对每个二进数 \(r = k/2^{n} \in [0, 1]\) ，设 \(A(r)\) 是由所有 \(x \in Y\) （满足 \(f(x) \leqslant r\) ）组成的集，并设 \(B(r) = A(r) \cup (X - Y)\) 。用归纳法 * （按 § 4 第 1 号定理 1 的证明方式） * 定义一个开集 \(U(r)\) （对每个二进数 \(r \in [0, 1]\) ），使得：只要 r 与 \(r'\) 是两个二进数且 \(r < r'\) ，就存在 X 的一致结构的一个围域 V，使得
\end{enumerate}

\[
\mathbf {V} (\mathbf {A} (r)) \subset \mathbf {U} (r ^ {\prime}), \quad \mathbf {V} (\mathbf {U} (r)) \subset \mathbf {U} (r ^ {\prime}), \quad \mathbf {V} (\mathbf {U} (r)) \subset \mathbf {B} (r ^ {\prime}).
\]

然后定义 \(\overline{f}(x)\) 为诸二进数 \(r\) ——满足 \(x \in \mathrm{U}(r).]\) 者——的下确界。

\ExerciseSection

\begin{enumerate}
\def\labelenumi{\arabic{enumi})}
\item
  设 X 是一个豪斯多夫一致空间，\((f_{t})\) 是 X 上定义 X 的一致结构的一族伪度量。设 \(X_{t}\) 是这样的一致空间相联系的度量空间：该一致空间由 X 赋以单个伪度量 \(f_{t}\) 所定义的结构而得（第 1 号）。证明 X 同构于乘积一致空间 \(\prod_{i} X_{i}\) 的一个子空间。
\item
  在一个连通的度量空间 X 中，若度量在 \(X \times X\) 上不是有界的，证明球面不可能是空的。
\end{enumerate}

¶ 3) a) 设 X 是一个紧度量空间。若每个开球在 X 中的闭包都是同中心同半径的闭球，则 X 中每个球都是连通集{[}若 x 与 y 是 X 的两点，S 是以 x 为中心、\(d(x, y)\) 为半径的闭球，证明对每个 \(\varepsilon > 0\) ，点所成的集 \(A_{y,\varepsilon}\) ——S 中可用一条含于 S 的 \(V_{\varepsilon}\) 链与 y 相连的那些点——（第 II 章，§ 4，第 4 号）含有满足 \(d(x, z) < d(x, y)\) 的点 z {]}。用反证法并利用魏尔斯特拉斯定理（第 IV 章，§ 6，第 1 号，定理 1），推出 \(x \in A_{y,\varepsilon}\) 对所有 \(\varepsilon > 0\) 成立。

\begin{enumerate}
\def\labelenumi{\alph{enumi})}
\setcounter{enumi}{1}
\tightlist
\item
  在空间 \(\mathbf{Y} = \mathbf{R}^2\) （赋以度量
\end{enumerate}

\[
d (x, y) = \max \left(\left| x _ {1} - y _ {1} \right|, \left| x _ {2} - y _ {2} \right|\right),
\]

）中，设 X 是由所有满足下列条件的点 \((x_{1}, x_{2})\) 组成的紧子空间：或者 \(x_{1} = 0\) 且 \(0 \leqslant x_{2} \leqslant 1\) ，或者 \(0 \leqslant x_{1} \leqslant 1\) 且 \(x_{2} = 0\) 。证明 X 中每个球都连通，但一个开球的闭包未必是同中心同半径的闭球。

\begin{enumerate}
\def\labelenumi{\arabic{enumi})}
\setcounter{enumi}{3}
\tightlist
\item
  度量空间 \(\mathbf{X}\) 称为超度量空间，如果它的度量 \(d\) 满足不等式
\end{enumerate}

\[
d (x, y) \leqslant \sup \left(d (x, z), d (y, z)\right)
\]

（对所有 \(x, y, z\) ∈ \(\mathbf{X}\) ；该不等式蕴含三角不等式）（参看 § 6，习题 2）。

\begin{enumerate}
\def\labelenumi{\alph{enumi})}
\item
  若 \(d(x, z) \neq d(y, z)\) ，证明 \(d(x, y) = \sup (d(x, z), d(y, z))\) 。
\item
  设 \(\mathrm{V}_{r}(x)\) 是以 x 为中心、r 为半径的开球。证明 \(\mathrm{V}_{r}(x)\) 在 X 中既开又闭（因而 X 是全不连通的），且对每个 \(y \in \mathrm{V}_{r}(x)\) 有 \(\mathrm{V}_{r}(y) = \mathrm{V}_{r}(x)\) 。
\item
  证明以 x 为中心、r 为半径的闭球 \(\mathrm{W}_{r}(x)\) 在 X 中既开又闭，且对每个 \(y \in \mathrm{W}_{r}(x)\) 有
\end{enumerate}

\[
\mathrm{W} _ {r} (y) = \mathrm{W} _ {r} (x).
\]

含于 \(W_{r}(x)\) 中的、半径为 r 的各个不同开球构成 \(W_{r}(x)\) 的一个分划，且这些球中任何两个之间的距离都等于 r 。d) 若 X 的两个球（开的或闭的）相交，则其中一个含于另一个之中。

\begin{enumerate}
\def\labelenumi{\alph{enumi})}
\setcounter{enumi}{4}
\item
  X 的点序列 \((x_{n})\) 是柯西序列，当且仅当 \(d(x_{n}, x_{n+1})\) 当 \(n\) 趋于无穷时趋于 0 。
\item
  若 X 是紧的，证明对每个 \(x_{0} \in X\) ，\(d(x_{0}, x)\) 在 X 中所取值的集合是 \([0, +\infty]\) 的一个（有限或无限的）可数子集，其所有点除 0 可能例外外都是孤立的{[}对 \(d(x_{0}, x)\) 所取的每个值 r，考虑 \(d(x_{0}, x)\) 在满足 \(d(x_{0}, x) < r\) 的点集上的上确界，以及它在满足 \(d(x_{0}, x) > r\) 的点集上的下确界{]}。
\end{enumerate}

\begin{enumerate}
\def\labelenumi{\arabic{enumi})}
\setcounter{enumi}{4}
\tightlist
\item
  设 X 是一个度量空间，d 是它的度量。对 X 的每一对点 \((x, y)\) ，令 \(d_{0}(x, y)\) 表示满足如下条件的实数 \(\alpha > 0\) 组成的集合的下确界：x 与 y 可以用一条 \(V_{\alpha}\) 链连接起来（第 II 章，§ 4，第 4 号）。证明 \(d_{0}\) 是 X 上的一个伪度量，且 X 上由伪度量 \(d_{0}\) 所定义的一致空间（第 1 号）相联系的度量空间是一个超度量空间（习题 4）。
\end{enumerate}

§ 6) 设 X 是一个度量空间。若 A 与 B 是 X 的两个非空子集，令

\[
\rho (\mathrm{A}, \mathrm{B}) = \sup _ {x \in \mathbf {A}} d (x, \mathrm{B}) \quad \text { and } \quad \sigma (\mathrm{A}, \mathrm{B}) = \sup \left(\rho (\mathrm{A}, \mathrm{B}), \rho (\mathrm{B}, \mathrm{A})\right);
\]

又令 \(\sigma(\emptyset, \emptyset) = 0\) ，\(\sigma(\emptyset, A) = \sigma(A, \emptyset) = +\infty\) （\(A\) 为 \(X\) 中任意非空子集）。证明 \(\sigma\) 是集合 \(\mathfrak{P}(X)\) （即 \(X\) 的全体子集组成的集合）上的一个伪度量，且由 \(\sigma\) 定义的一致结构与按第 II 章，§ 1，习题 5 的程序由 \(X\) 的一致结构构造出的一致结构相重合。

假设 X 是有界的。此时 \(\sigma\) 是集合 \(\mathfrak{F}(\mathbf{X})\) （X 的非空闭子集全体）上的一个度量。此外若 X 还是完备的，证明 \(\mathfrak{F}(\mathbf{X})\) 是一个完备度量空间{[}设 \(\Phi\) 是 \(\mathfrak{F}(\mathbf{X})\) 上的一个柯西滤子；对每个集 \(\mathcal{X} \in \Phi\) ，令 \(S(\mathcal{X})\) 表示属于 \(\mathcal{X}\) 的那些 X 的子集 A 的并集；证明诸集 \(S(\mathcal{X})\) 当 \(\mathcal{X}\) 跑遍 \(\Phi\) 时构成 X 上的一个滤子基，该滤子基具有一个非空的聚点集 C，且 \(\Phi\) 收敛于 C{]}。

\begin{enumerate}
\def\labelenumi{\arabic{enumi})}
\setcounter{enumi}{6}
\item
  设 X 是一个无限离散空间。证明 X 上的有限分划的一致结构（第 II 章，§ 2，第 2 号）——它与 X 的拓扑相容——不是可度量化的（否则存在 X 的有限分划的序列 \((\mathfrak{F}_n)\) ，使得 X 的每个有限分划都由一些集组成，其中每个集是属于诸分划 \(\mathfrak{F}_n\) 之一的集的并；由此将推出 X 的有限分划组成的集合是可数的）。
\item
  设 \((\mathbf{X}_{t})\) 是一个不可数的豪斯多夫拓扑空间族，其中每个空间至少有两个不同的点。证明在乘积空间 \(\prod_{i} X_{i}\) 中，任何点都没有可数的基本邻域系。
\item
  设 X 是一个拓扑空间，它的每个点都有可数的基本邻域系。
\end{enumerate}

\begin{enumerate}
\def\labelenumi{\alph{enumi})}
\item
  只要 X 中每个收敛序列都只有一个极限，X 就是豪斯多夫的。
\item
  若 a 是 X 的点组成的无限序列 \((x_{n})\) 的一个聚点，则存在 \((x_{n})\) 的一个无限子序列收敛于 a。
\item
  设 A 是 X 的一个非空子集，\(x_{0}\) 是 \(\overline{A}\) 的一个点，f 是 A 到一个豪斯多夫拓扑空间 \(X'\) 中的映射。则点 \(a \in X'\) 是 f 在点 \(x_{0}\) 处（关于 A）的极限，只要对 A 中点的每个序列 \((x_{n})\) （它收敛于 \(x_{0}\) ），序列 \((f(x_{n}))\) 都收敛于 a。
\item
  沿用 c) 的记号，再设 \(X'\) 的每个点都有可数的基本邻域系。若 \(a \in X'\) 是 f 在 \(x_{0}\) 处关于 A 的一个聚点，则存在 A 的点组成的序列 \((x_{n})\) ，它收敛于 \(x_{0}\) ，且使序列 \((f(x_{n}))\) 收敛于 a。
\end{enumerate}

¶ 10) 设 X 是一个紧空间。若存在 X×X 上的一个连续实值函数 f，使得关系 \(f(x,y)=0\) 等价于 x=y，则 X 是可度量化的。（证明若 V 跑遍 R 中 0 的一个基本邻域系，则诸集 \(\overline{f}^{1}(V)\) 构成 X×X 中对角线 \(\Delta\) 的一个基本邻域系。）

¶ 11) 设 X 是一个紧度量空间，d 是它的度量。证明：若 f 是 X 到 X 中的映射，使得对 X 的每一对点 x、y 有 \(d(f(x), f(y)) \geqslant d(x, y)\) ，则 f 是 X 到其自身上的一个等距映射。{[}设 a、b 是 X 的任意两点，令 \(f^n = f^{n-1} \circ f\) ，\(a_n = f^n(a)\) ，\(b_n = f^n(b)\) ；证明对每个 \(\varepsilon > 0\) ，存在指标 k 使得 \(d(a, a_k) \leqslant \varepsilon\) 且 \(d(b, b_k) \leqslant \varepsilon\) （选取 \((a_n)\) 与 \((b_n)\) 的适当子序列）；由此推出 \(d(a_1, b_1) = d(a, b)\) ，且 \(f(\mathbf{X})\) 在 X 中稠密。{]}

¶ 12) 设 X 是一个紧可度量化空间。

\begin{enumerate}
\def\labelenumi{\alph{enumi})}
\item
  证明存在康托尔三分集 K（第 IV 章，§ 2，第 5 号）到 X 上的一个连续映射（参看第 IV 章，§ 8，习题 11）。
\item
  若此外 X 还是全不连通的且没有孤立点，则 X 同胚于 K。{[}按第 IV 章，§ 8，习题 12 的方式论证，利用如下事实：点 \(x \in \mathbf{X}\) 的每个邻域都包含 \(x\) 的一个既开又闭的邻域（第 II 章，§ 4，第 4 号，命题 6 的推论）{]}。
\end{enumerate}

¶ 13) a) 在实数集 R 上，考虑如下定义的拓扑 C：对每个 y \textgreater{} 0，\(\mathrm{U}_{y}(x)\) 表示诸区间的并 \([x, x + y[ \text{ 和 } ] - x - y, -x[; C \text{ 是满足下述条件的拓扑： } \mathrm{U}_{y}(x) \text{ 构成邻域基本系，其中心为 } x \text{ 当 } y \text{ 遍历实数集 } > 0. \text{ 令 X 为在区间 } [-1, +1] \text{ 上赋予由 C 诱导的拓扑所得到的空间。证明 X 是紧的。（考虑 X 上的一个滤子 F；若 } x \in X \text{ 是 F 在实直线拓扑下的聚点，证明 } x \text{ 或 } -x \text{ 是 F 在 C 下的聚点.)}\)

\begin{enumerate}
\def\labelenumi{\alph{enumi})}
\setcounter{enumi}{1}
\item
  X 的每个点都有可数的基本邻域系，且 X 有可数稠密子集，但 X 不是可度量化的（证明它的拓扑没有可数基）。
\item
  设 A 是 X 的一个开子集。证明 A 是如下诸集之并：开区间组成的一个可数族 \((I_{n})\) （这些区间含于 {[}0, 1{]} ）、诸区间 \(-I_{n}\) 、诸区间 \(I_{n}\) 与 \(-I_{n}\) 的左端点组成的集合的一个子集，以及可能还有点 +1。由此证明 A 是 X 的闭子集组成的一个可数族之并。
\item
  设 Y 是集合 \(I \times \{1, 2\}\) ，其中 I 是 R 的区间 {[}0, 1{]} ；令 \(X_{i}\) 表示 \(I \times \{i\} (i = 1, 2)\) ，令 f 表示双射 \((x, 1) \to (x, 2)\) （它是 \(X_{1}\) 到 \(X_{2}\) 上的映射）。对每个 \(x \in I\) ，令 \(\mathfrak{B}((x, 2))\) 表示仅由 \(\{(x, 2)\}\) 组成的子集所成的集合，令 \(\mathfrak{B}((x, 1))\) 表示形如 \(V_{1} \cup (f(V_{1}) - \{x, 2\})\) 的子集所成的集合，其中 \(V_{1} = V \times \{1\}\) 而 V 跑遍 I 中 x 的一个基本邻域系。证明对每个 \(y \in Y\) ，\(\mathfrak{B}(y)\) 是 Y 上某个拓扑的、y 的基本邻域系。赋以此拓扑后，Y 是紧的，且 Y 的每个点都有可数的基本邻域系，但 Y 没有可数稠密子集。Y 的每个含于 \(X_{2}\) 中的闭子集都是有限的；由此证明 Y 的紧子集 \(X_{1}\) 没有可数的基本邻域系。若 R 是 Y 上的等价关系，其等价类是集合 \(X_{1}\) 以及 \(Y - X_{1}\) 的诸点，则 R 是闭的，且模 R 的每个等价类都是紧的，但 Y/R 中有一个点没有可数的基本邻域系{[}参看 § 4，习题 24 a){]}。
\end{enumerate}

\begin{enumerate}
\def\labelenumi{\arabic{enumi})}
\setcounter{enumi}{13}
\tightlist
\item
  \begin{enumerate}
  \def\labelenumii{\alph{enumii})}
  \tightlist
  \item
    设 X 是一个可达的拓扑空间（第 I 章，§ 8，习题 1）。证明下列性质等价：
  \end{enumerate}
\end{enumerate}

α) X 的每个点序列都有聚点。

\(\beta)\) \(\mathbf{X}\) 的任何无限离散子空间都不是闭的。

\(\gamma)\) \(\mathbf{X}\) 的每个可数开覆盖都包含 \(\mathbf{X}\) 的一个有限开覆盖。

δ) 对 X 的任何无限开覆盖 R，存在 X 的一个开覆盖 S ⊂ R，它异于 R。

拓扑空间 X 称为可数紧的，如果它是豪斯多夫的且具有上述性质。

\begin{enumerate}
\def\labelenumi{\alph{enumi})}
\setcounter{enumi}{1}
\item
  可数紧空间中的点序列是收敛的，当且仅当它有唯一的聚点。
\item
  可数紧空间的每个闭子空间都是可数紧的。反之，若 X 是豪斯多夫的且 X 的每个点都有可数的基本邻域系，则 X 的每个可数紧子空间在 X 中是闭的。
\item
  设 \(f\) 是可数紧空间 \(\mathbf{X}\) 到豪斯多夫空间 \(\mathbf{X}'\) 中的一个连续映射。则 \(f(\mathbf{X})\) 是 \(\mathbf{X}'\) 的一个可数紧子空间。
\item
  设 X 是一个可数紧空间，它的每个点都有可数的基本邻域系。则 X 的每个点序列都有收敛子序列。
\end{enumerate}

\(f\) ) 设 \((\mathbf{X}_n)\) 是拓扑空间组成的一个可数序列，其中每个空间的每个点都有可数的基本邻域系。

此时乘积空间 \(\prod_{n=1}^{\infty} \mathbf{X}_n\) 是可数紧的，当且仅当诸空间 \(\mathbf{X}_n\) 中的每一个都是可数紧的。{[}利用 e) 及第 I 章，§ 6，习题 16。{]}

\begin{enumerate}
\def\labelenumi{\alph{enumi})}
\setcounter{enumi}{6}
\item
  每个点都有可数基本邻域系的可数紧空间是正则的。
\item
  若一个可数紧空间具有可数的开集基，则它是紧的，从而是可度量化的。
\item
  可数紧空间上的每个下半连续实值函数都达到它的下确界。特别地，每个可数紧空间都是伪紧的（§ 1，习题 21）{[}参看 § 4，习题 26 b){]}。
\item
  证明性质 \(\delta\) （见 \(a\) ）蕴含下列性质：
\end{enumerate}

ζ) X 的每个逐点有限（§ 4，第 3 号）开覆盖都包含 X 的一个有限开覆盖。（用反证法论证。）反之，满足 ζ) 的正则空间是可数紧的{[}证明此时 ζ) 蕴含 β){]}。

¶ 15) 设 \(\mathbf{X} = [a, b[\) 是第 I 章，§ 9，习题 12 中定义的局部紧空间。

\begin{enumerate}
\def\labelenumi{\alph{enumi})}
\item
  X 的一个子集是相对紧的，当且仅当它是有界的。由此证明 X 是可数紧的，从而（命题 15）不是可度量化的（注意 X 的每个可数子集都是有界的）。
\item
  证明 X 的每个点都有一个可度量化邻域（利用命题 16）。
\item
  若 A 与 B 是 X 中两个非紧的闭集，则它们的交不是空的（构造一个递增序列，使其偶数编号的项是 A 的点，奇数编号的项是 B 的点）。由此推出，非紧闭集的每个邻域都是一个相对紧集的余集。若 A 是 X 的非孤立点组成的集合，证明 A 是闭的，且不是任何可数开集族的交。
\item
  令 \(\mathbf{X}'\) 表示赋有如下拓扑的区间 \([a, b]\) ：对每个 \(x \in \mathbf{X}\) ，诸区间 \(]y, x]\) （其中 \(y\) 跑遍由满足 \(< x\) 的元素组成的集合）构成 \(x\) 的一个基本邻域系；\(b\) 的一个基本邻域系由诸集 \(\mathbf{V}_x\) 组成，其中对每个 \(x \in \mathbf{X}\) ，\(\mathbf{V}_x\) 表示 \(\{b\}\) 与 \([x, b]\) 中具有前元的点（即在 \(\mathbf{X}\) 中孤立的那些点）组成的集合之并。证明 \(\mathbf{X}'\) 是可数紧的但不是正则的{[}参看习题 14 g){]}，且子空间 \(\mathbf{X}\) （它含于 \(\mathbf{X}'\) ）是可数紧的但不是闭的{[}参看习题 14 c){]}。
\end{enumerate}

\begin{enumerate}
\def\labelenumi{\arabic{enumi})}
\setcounter{enumi}{15}
\tightlist
\item
  设 X 与 Y 是两个可数紧空间。证明在下列两种情形的每一种中，乘积 \(X \times Y\) 都是可数紧的：(i) X、Y 之一是紧的；(ii) X、Y 之一使每个点都有可数的基本邻域系。{[}在情形 (i)，设 Y 是紧的，考虑一个递增的开集序列 \((\mathbf{G}_{n})\) ，其并为 \(X \times Y\) ；对每个 n，令 \(H_{n}\) 表示所有满足如下条件的 \(x \in X\) 组成的集合：\(\{x\} \times Y \subset G_{n}\) ；证明 \(H_{n}\) 是开的，且 \(X = \bigcup_{n} H_{n} \cdot\) {]}
\end{enumerate}

¶ 17) 设 βN 是离散空间 N 的斯通–切赫紧化（§ 1，习题 7）。

\begin{enumerate}
\def\labelenumi{\alph{enumi})}
\item
  证明存在 \(\beta N\) 的两个可数紧子空间 A、B，使得 \(A \cap B = N\) 且 \(A \cup B = \beta N\) 。{[}令 \(\aleph_{\alpha} = 2^{2^{\operatorname{Card}(N)}}\) ；根据 § 1，习题 12 b)，存在一个双射 \(\xi \to S_{\xi}\) ，它把序数 \(\omega_{\alpha}\) （《集合论》R, 第 III 章，§ 6，习题 10）映到 X 的可数无限子集组成的集合上。用超限归纳法定义两个单射映射 \(\xi \to x_{\xi}, \xi \to y_{\xi}\) ，它们把 \(\omega_{\alpha}\) 映入 \(\beta N - N\) ，使得 \(x_{\xi} \in \overline{S}_{\xi}, y_{\xi} \in \overline{S}_{\xi}\) 对每个 \(\xi\) 成立，且使得若 P（相应地 Q）是诸 \(x_{\xi}\) （相应地诸 \(y_{\xi}\) ）组成的集合，则有 \(P \cap Q = \varnothing\) ，\(P \cup Q = \beta N - N\) ；为此，利用 § 1，习题 12 b) 与 c)。然后证明 A = P ∪ N 与 B = Q ∪ N 满足问题的条件。{]}
\item
  证明乘积空间 \(A \times B\) 不是可数紧的（注意 \(A \times B\) 与 \(\beta N \times \beta N\) 的对角线的交是 \(A \times B\) 的一个无限离散闭子空间）。
\end{enumerate}

¶ 18) 设 X 是一个度量空间，使得对每个 \(x \in X\) ，存在一个以 \(x\) 为中心的开球，它作为 X 的子空间考虑时具有可数基。令 \(r_x\) 表示具有该性质的、以 \(x\) 为中心的球的半径所组成的集合的上确界。

\begin{enumerate}
\def\labelenumi{\alph{enumi})}
\item
  利用佐恩引理证明：存在 X 中两两不相交的开球组成的一个极大族 \((\mathbf{B}_{\alpha})\) ，使得若 \(x_{\alpha}\) 是 \(B_{\alpha}\) 的中心，则 \(B_{\alpha}\) 的半径 \(<r_{x_{\alpha}}\) 。证明诸 \(B_{\alpha}\) 的并在 X 中稠密，从而存在 X 的一个稠密子集 M，使得以任意点 \(x \in X\) 为中心、半径 \(<r_{x}\) 的每个开球只含 M 的可数无穷多个点（注意这样一个球只能与可数无穷多个两两不相交的开集相交，并利用命题 12）。
\item
  对每个点 \(x \in \mathbf{M}\) ，令 \(\mathbf{S}_x\) 表示以 \(x\) 为中心、\(\frac{1}{3} r_x\) 为半径的开球。证明诸 \(\mathbf{S}_x\) 覆盖 \(\mathbf{X}\) ，且 \(\mathbf{X}\) 的任意点只属于可数无穷多个 \(\mathbf{S}_x\) {[}注意若 \(y \in \mathbf{S}_x\) ，则有 \(d(x, y) \leqslant \frac{1}{2} r_y\) {]}。
\item
  设 R 是 X 的点 \(x, y\) 之间的如下等价关系：存在 M 的点组成的序列 \((z_i)_{1 \leqslant i \leqslant n}\) ，使得 \(x \in S_{z_i}\) ，\(y \in S_{z_n}\) ，且 \(S_{z_i}\) 与 \(S_{z_{i+1}}\) 相交（对 \(1 \leqslant i \leqslant n - 1\) ）。证明（模 R 的）等价类是 X 的一些子空间，它们在 X 中既开又闭，且具有可数基；换言之，X 是具有可数基的一些度量空间的拓扑和（第 I 章，§ 2，第 4 号）。
\end{enumerate}

¶ 19) 在度量空间中，每个相对紧集都是有界的。若 X 是一个拓扑空间，证明 X 的下列条件等价：α) 存在 X 上的一个度量，它与 X 的拓扑相容，使得 X 的每个（关于该度量的）有界子集都是相对紧的；β) X 是局部紧的且具有可数基。{[}为证 α) → β)，证明若每个有界集都相对紧，则 X 是局部紧的且是 σ-紧的。为证 β) → α)，注意由命题 12 的推论，X 是可度量化的；若 d 是 X 上与 X 的拓扑相容的一个度量，且 f 是 § 1，习题 15 d) 中定义的函数，则取 X 上由两个伪度量 d(x, y) 与 \textbar f(x) --- f(y)\textbar 所定义的一致结构。{]}

¶ 20) a) 给出一个例子：在具有可数基的可度量化局部紧空间 X 上的一个闭等价关系 R，使得 X/R 是仿紧的，但 X/R 中存在一个点没有可数的基本邻域系{[}参看第 I 章，§ 10，习题 17 及第 IX 章，§ 4，习题 24 a){]}。

\begin{enumerate}
\def\labelenumi{\alph{enumi})}
\setcounter{enumi}{1}
\tightlist
\item
  设 X 是一个可度量化空间，R 是 X 上的一个闭等价关系。证明：若 X/R 的每个点都有可数的基本邻域系，则每个（模 R 的）等价类在 X 中具有紧的边界。（证明：若结论不真，则存在 X 中的一个序列，其点两两不同、没有聚点，而它在 X/R 中的像是一个序列，收敛到与该序列所有点都不同的一个点。）对每个 \(z \in X/R\) ，令 \(C_{z}\) 为 z 在 X 中的逆像，且令 \(F_{z}\) 为 \(C_{z}\) 的边界（若该边界非空）；若 \(C_{z}\) 既开又闭，令 \(F_{z}\) 为由 \(C_{z}\) 的单独一个点组成的任一子集。令 Y 为当 z 跑遍 X/R 时诸集 \(F_{z}\) 的并。证明 \(Y/R_{y}\) 同胚于 X/R。
\end{enumerate}

¶ 21) a) 设 X 是一个可度量化空间，d 是与 X 的拓扑相容的一个有界度量，σ 是 d 在非空闭子集组成的集合 ℱ(X) 上对应的度量（习题 6）；R 是 X 上的一个既开又闭的等价关系。证明 σ 在 X/R 上的限制是与商拓扑相容的一个度量。{[}利用习题 20 b) 证明模 R 的每个等价类都是开的且紧的。然后证明：若 X/R 中的序列 (zn) 趋于一个点 a，且 φ: X → X/R 是典范映射，则有 σ(φ̄¹(a), φ̄¹(zₙ)) → 0；利用 R 是开的这一事实，用反证法论证。{]}

\begin{enumerate}
\def\labelenumi{\alph{enumi})}
\setcounter{enumi}{1}
\tightlist
\item
  在 R 的紧区间 I = {[}0, 1{]} 上，设 R 是这样的等价关系：其等价类为康托尔集 K（第 VI 章，§ 2，第 5 号）中除 K 的邻接区间的端点以外的点，以及 K 的邻接区间的闭包。证明商空间 I/R 同胚于 I{[}第 IV 章，§ 8，习题 16 b){]}，但度量 \(\sigma\) 与 I/R 的拓扑不相容。
\end{enumerate}

\begin{enumerate}
\def\labelenumi{\arabic{enumi})}
\setcounter{enumi}{21}
\tightlist
\item
  设 X 是具有可数基 \((\mathbf{U}_{n})\) 的一个豪斯多夫拓扑空间，R 是 X 上的一个豪斯多夫等价关系，使得 X/R 的每个点都有可数的基本邻域系。证明 X/R 的拓扑具有可数基。{[}设 \(\varphi: X \to X/R\) 是典范映射，证明诸集 \(\varphi(\mathbf{U}_{n})\) 的有限并的内部构成 X/R 的拓扑的一个基：若 V 是点 \(z \in X/R\) 的一个邻域，且 \((\mathbf{W}_{k})\) 是属于基 \((\mathbf{U}_{n})\) 、含于 \(\overline{\varphi}^{1}(\mathbf{V})\) 且覆盖 \(\overline{\varphi}^{1}(z)\) 的集组成的一个序列，证明存在有限多个指标 k，使诸 \(\varphi(\mathbf{W}_{k})\) 的并是 z 的一个邻域。用反证法：若此命题不真，则可构造 X/R 的两两不同的点组成的序列 \((y_{n})\) ，它趋于 z，且使 \(y_{n}\) 不属于诸 \(\varphi(\mathbf{W}_{k})\) （\(k \leqslant n\) ）的并；然后证明诸 \(\overline{\varphi}^{1}(y_{n})\) 的并在 X 中是闭的。{]}
\end{enumerate}

证明若假设 R 是闭的且模 R 的每个等价类都是紧的，则同样的结论成立（方法类似）。

¶ 23) 拓扑空间称为次可度量化的，如果它的拓扑比某个可度量化空间的拓扑更细。次可度量化空间是豪斯多夫的，但未必是正则的（第 I 章，§ 8，习题 20）。

\begin{enumerate}
\def\labelenumi{\alph{enumi})}
\item
  完全正则空间 X 是次可度量化的，当且仅当存在 X 上的一个一致结构，它与 X 的拓扑相容，且由一个伪度量族 \(\Phi\) 定义，该族至少包含一个度量 d（参看 § 1，习题 3）。设 \(\hat{X}\) 是 X 关于这个一致结构的完备化；度量 d 可扩张为一个伪度量 \(\overline{d}\) （在 \(\hat{X}\) 上）；证明对每个 \(x \in \hat{X}\) ，至多存在一个点 \(y \in X\) 使得 \(\overline{d}(x, y) = 0\) 。
\item
  证明若 X 是一个完全正则的次可度量化空间，则存在与 X 的拓扑相容的一个一致结构，X 关于它是完备的{[}利用 a) 及 § 1，习题 16 c){]}。由此证明若此外 X 还是伪紧的（§ 1，习题 21），则 X 是紧的。
\item
  证明局部紧、仿紧、次可度量化的空间 X 是可度量化的（利用第 I 章，§ 9，第 10 号，定理 5 化归为 X 是 σ-紧的情形；然后利用命题 16 的推论）{[}参看 § 5，习题 15{]}。
\end{enumerate}

¶ 24) a) 设 X 是一个完备度量空间，A 是 X 的一个子集，它是一个可数开集族的交。证明 A 上存在一个度量，A 关于它是一个完备度量空间，且它在 A 上既定义由 X 的拓扑诱导的拓扑，又定义一个比 X 的一致结构所诱导的一致结构更细的一致结构{[}参看 § 1，习题 16 b) 与 c){]}。

\begin{enumerate}
\def\labelenumi{\alph{enumi})}
\setcounter{enumi}{1}
\tightlist
\item
  反之，设 X 是一个可度量化空间，A 是 X 的一个子集，使得 A 上存在一个度量 d，它与由 X 的拓扑诱导的拓扑相容，且 A 关于它是一个完备度量空间。证明 A 是 X 中可数个开集的交。{[}对每个整数 n \textgreater{} 0，考虑集合 \(G_{n}\) ，它由具有如下性质的点 \(x \in \overline{A}\) 组成：x 有一个开邻域 U，使 \(A \cap U\) （关于 d）的直径 \(\leqslant r/n\) 。{]}
\end{enumerate}

¶ 25) 设 X 是一个可度量化空间，βX 是它的斯通–切赫紧化（§ 1，习题 7），d 是与 X 的拓扑相容的一个有界度量。对每个 x ∈ X，令 f\_x(y) 表示把 y → d(x, y) 连续地扩张到 βX 上所得的实值函数。

\begin{enumerate}
\def\labelenumi{\alph{enumi})}
\item
  证明 \(f_{x}(z) + f_{y}(z) \geqslant d(x, y)\) 对所有 \(x\) 、\(y\) ∈ \(\mathbf{X}\) 及 \(z\) ∈ \(\beta \mathbf{X}\) 成立。对每个 \(y \in \beta \mathbf{X}\) （异于 \(x \in \beta \mathbf{X}\) ），证明 \(f_{x}(y) > 0\) 。b) 证明若 \(\mathbf{X}\) 关于度量 \(d\) 是完备的，则 \(\mathbf{X}\) 是 \(\beta \mathbf{X}\) 中一个开集序列的交。{[}考虑集合 \(G_{n}\) ，它由所有满足下列条件的 \(y \in \beta \mathbf{X}\) 组成：\(f_{x}(y) < 1 / n\) 对至少一个点 \(x \in \mathbf{X}\) 成立，并利用 a) 证明 \(\mathbf{X} = \bigcap_{n=1}^{\infty} G_{n}\) 。{]}
\item
  反之，证明若 X 是一个开集序列 \(G_{n}\) （含于 \(\beta X\) ）的交，则 X 上存在一个度量 \(d'\) ，它与 X 的拓扑相容，且 X 关于它是完备的。{[}只需考虑 \(X \neq \beta X\) 的情形，注意 \(\beta X - X\) 是紧集组成的一个族 \(F_{n} = \beta X - G_{n}\) 的并；对每个 n，令 \(g_{n}(x) = \inf_{y \in F_{n}} f_{x}(y)\) ，并利用 a) 证明 \(g_{n}(x) > 0\) 对所有 \(x \in X\) 成立，且 \(g_{n}\) 在 X 上连续。证明的结尾按 § 1，习题 16 b) 的方式进行。{]}
\end{enumerate}

\ExerciseSection

\begin{enumerate}
\def\labelenumi{\arabic{enumi})}
\tightlist
\item
  伪度量 \(f\) （在群 \(G\) 上，乘法记号）称为左不变的（相应地，右不变的），如果 \(f(zx, zy) = f(x, y)\) {[}相应地，\(f(xz, yz) = f(x, y)\) {]} 对所有 \(x, y, z\) ∈ \(G\) 成立。
\end{enumerate}

\begin{enumerate}
\def\labelenumi{\alph{enumi})}
\tightlist
\item
  若 \(f\) 是一个左不变伪度量，则实值函数 \(g(x) = f(e, x)\) （在 \(\mathbf{G}\) 上；\(e\) 是 \(\mathbf{G}\) 的单位元）满足下列条件：
\end{enumerate}

\begin{enumerate}
\def\labelenumi{(\roman{enumi})}
\item
  \(g(x)\geqslant \mathrm{0}\) 对所有 \(x\in \mathbf{G}\) 成立，且 \(g(e) = \mathrm{0};\)
\item
  \(g(x^{-1}) = g(x)\) ；
\item
  \(g(xy) \leqslant g(x) + g(y)\) 。
\end{enumerate}

反之，若 \(g\) 是 \(G\) 上满足这些条件的任一实值函数，则 \(f(x, y) = g(x^{-1}y)\) 是 \(G\) 上的一个左不变伪度量。\(b)\) 拓扑 \(\mathfrak{G}\) ——由左不变伪度量组成的一个饱和族 \((f_t)\) （在群 \(G\) 上）所定义——与 \(G\) 的群结构相容，当且仅当对每个 \(a \in G\) 、每个指标 \(t\) 与每个实数 \(\alpha > 0\) ，存在一个指标 \(x\) 与一个实数 \(\beta > 0\) ，使得关系 \(f_x(e, x) \leqslant \beta\) 蕴含 \(f_t(e, axa^{-1}) \leqslant \alpha\) 。若此条件满足，则 \(G\) 上由族 \(f_t\) 定义的一致结构，与把 \(G\) 赋以拓扑 \(\mathfrak{G}\) 所得的拓扑群的左一致结构相同。

\begin{enumerate}
\def\labelenumi{\alph{enumi})}
\setcounter{enumi}{2}
\tightlist
\item
  在每个拓扑群 G 上，存在左不变伪度量组成的一个族，使得由该族定义的一致结构与 G 的左一致结构相同。
\end{enumerate}

\begin{enumerate}
\def\labelenumi{\arabic{enumi})}
\setcounter{enumi}{1}
\item
  设 \(\mathbf{G}\) 是一个左、右一致结构相重合的拓扑群。证明这个唯一的一致结构可以由一个同时左不变且右不变的伪度量族定义（利用第 III 章，§ 3，习题 3，证明：若 \(\mathbf{V}\) 是 \(e\) （\(\mathbf{G}\) 的单位元）的任一邻域，则 \(\mathbf{V}_0 = \bigcap_{x\in \mathbf{G}}x\mathbf{V}x^{-1}\) 是 \(e\) 的一个邻域）。
\item
  设 G 是一个拓扑群，\((f_{t})\) 是 G 上的定义 G 的左一致结构的一个左不变伪度量饱和族，且令 \(g_{t}(x)=f_{t}(e,x)\) 。设 H 是 G 的一个闭正规子群，对每个陪集 \(\dot{x}\in G/H\) ，令 \(h_{t}(\dot{x})=\inf_{x\in \dot{x}}g_{t}(x)\) ；证明：若 \(\bar{f}_{t}(\dot{x},\dot{y})=h_{t}(\dot{x}^{-1}\dot{y})\) ，则诸 \(\bar{f}_{t}\) 构成 G/H 上的一个左不变伪度量族，它定义 G/H 的左一致结构（按第 1 号的注记 2 的方式论证）。
\item
  证明赋值除环的拓扑是局部逆有界的（第 III 章，§ 6，习题 22）。
\end{enumerate}

¶ 5) 设 \(\omega\) 是有理数域 \(\mathbf{Q}\) 上的一个绝对值。a) 证明 \(\omega\) 由它在素数处的值唯一确定。b) 若存在一个素数 \(p\) 使 \(\omega(p) \leqslant 1\) ，证明 \(\omega(q) \leqslant 1\) 对每个其他素数 \(q\) 成立{[}给出 \(\omega(q^n)\) 的一个上界，办法是把 \(q^n\) 写成以 \(p\) 为基的尺度，然后令 \(n \to \infty\) {]}。

\begin{enumerate}
\def\labelenumi{\alph{enumi})}
\setcounter{enumi}{2}
\tightlist
\item
  若存在一个素数 p 使 \(\omega(p) < 1\) ，证明 \(\omega(q) = 1\) 对每个其他素数 q 成立{[}利用如下事实证明不可能有 \(\omega(q) < 1\) ：对每个整数 n \textgreater{} 0，存在两个有理整数 r 与 s 使得 \(1 = rp^{n} + sq^{n}\) {]} 。由此证明 \(\omega\) 此时是一个等价于 Q 上的 p 进绝对值的绝对值。d) 若对每个素数 p 有 \(\omega(p) > 1\) ，证明对任意两个素数 p 与 q，有
\end{enumerate}

\[
\frac {\log \omega (p)}{\log p} = \frac {\log \omega (q)}{\log q}
\]

{[}方法与 \(b\) 中相同{]}。由此证明在此情形 \(\omega(x) = |x|^{\rho}\) 对某个 \(\rho \leqslant 1\) 成立。

\begin{enumerate}
\def\labelenumi{\arabic{enumi})}
\setcounter{enumi}{5}
\item
  设 E 是除环 K 上的一个左向量空间，具有可数基 \((a_{n})\) 。设 \((r_{n})\) 是一个递减的、趋于 0 的数组成的序列 \textgreater0。对 E 的每个 \(x = \sum_{k} t_{k} a_{k} \neq 0\) ，令 \(\|x\| = r_{h}\) ，其中 h 是使 \(t_{k} \neq 0\) 的诸指标 k 中最小者；并令 \(\|0\| = 0\) 。证明 \(\|x - y\|\) 是 E 的加法群上的一个不变度量，且它在 E 上定义的拓扑与所选取的序列 \((r_{n})\) （递减且趋于 0）无关。由此证明：若把范数与赋范空间的定义（定义 5 与定义 6）推广到标量上的绝对值是平凡绝对值的情形，则命题 7 与定理 1 不再成立。
\item
  设 E 是赋值除环上的一个赋范空间。证明：若 E 中每个绝对可和族都在 E 中可和，则 E 是完备的。{[}设 \((x_{n})\) 是 E 中的一个柯西序列，考虑一个子序列 \((x_{n_{k}})\) （取自 \((x_{n})\) ），使得以 \(x_{n_{k+1}} - x_{n}\) 为通项的级数绝对收敛。{]}
\item
  给出 p 进数域 \(Q_{p}\) 中一个可和但不绝对可和的族的例子。{[}证明族 \((x_{t})_{t\in\mathbf{I}}\) （\(Q_{p}\) 的点组成的）是可和的，当且仅当 \(\lim x_{t}=0\) 关于 I 的有限子集的余集组成的滤子成立。{]}
\item
  映射 \(w\) （从环 \(\mathbf{A}\) 到 \(\mathbf{R}_+\) 中）称为 \(\mathbf{A}\) 上的半绝对值，如果它满足条件：(i) \(w(\mathrm{0}) = \mathrm{0}\) ；(ii) \(w(x - y) \leqslant w(x) + w(y)\) ；
\end{enumerate}

\begin{enumerate}
\def\labelenumi{(\roman{enumi})}
\setcounter{enumi}{2}
\tightlist
\item
  \(w(xy) \leqslant w(x)w(y)\) 。半绝对值 w 称为豪斯多夫的，如果 \(w(x) = 0\) 蕴含 x = 0。若 w 是 A 上的一个豪斯多夫半绝对值，则 \(w(x - y)\) 是 A 的加法群上的一个不变度量，从而定义一个与 A 的加法群结构相容的拓扑。把赋范代数的主要性质，特别是命题 13，推广到赋有半绝对值的环上。
\end{enumerate}

若 \(w\) 是 \(A\) 上一个非豪斯多夫的半绝对值，则集合 \(\overline{w}(0)\) 是一个双边理想 \(\mathfrak{a}\) （在 \(A\) 中）；若给 \(A\) 赋以由伪度量 \(w(x - y)\) 定义的拓扑，则该拓扑与 \(A\) 的环结构相容，且与 \(A\) 相伴的豪斯多夫空间是商环 \(A / \mathfrak{a}\) ，在它上面，函数 \(\overline{w}(\overline{x})\) ——对每个 \(\overline{x} \mod \mathfrak{a}\) ，它等于 \(w(x)\) （对所有 \(x \in \overline{x}\) ）的公共值——是一个称为与 \(w\) 相伴的豪斯多夫半绝对值；由 \(\overline{w}\) 定义的拓扑是商去 \(\mathfrak{a}\) 之后的 \(A\) 的拓扑。

¶ 10) a) 环 A 上的两个半绝对值 \(w_1, w_2\) 称为等价的，如果伪度量 \(w_1(x - y), w_2(x - y)\) 是等价的。证明：若 \(w\) 是 A 上的一个半绝对值，则 \(aw\) 与 \(w^{1/a}\) 是等价于 \(w\) 的半绝对值（对每个实数 \(a \geqslant 1\) ）。

\begin{enumerate}
\def\labelenumi{\alph{enumi})}
\setcounter{enumi}{1}
\tightlist
\item
  若 \(w_{i}\) （ \(1 \leqslant i \leqslant n\) ）是环 A 上的半绝对值，则函数 \(w = \sum_{i} w_{i}\) 与 \(w' = \sup_{i} w_{i}\) 是 A 上两个等价的半绝对值。若
\end{enumerate}

\[
\mathfrak {a} _ {i} = \overline {{{w}}} _ {i} ^ {1} (\mathrm{0}) \quad \text { and } \quad \mathfrak {a} = \overline {{{w}}} ^ {1} (\mathrm{0}),
\]

则有 \(\mathfrak{a} = \bigcap_{i} \mathfrak{a}_{i}\) 。若 \(A_i\) 表示商环 \(A/a_i\) （赋以与 \(w_i\) 相伴的豪斯多夫半绝对值）的完备化，证明 \(A/a_i\) （赋以与 \(w\) 相伴的豪斯多夫半绝对值）同构于乘积环 \(\prod_{i} A_i\) 的一个子环；证明（假设 \(A\) 有单位元 \(1\) ）这个环在 \(\prod_{i} A_i\) 中稠密，当且仅当诸 \(w_i\) 满足下列条件：对每个 \(\varepsilon > 0\) 及每个指标 \(i\) ，存在一个元素 \(x_i \in A\) ，使得 \(w_i(1 - x_i) \leqslant \varepsilon\) ，且 \(w_k(x_i) \leqslant \varepsilon\) 对每个指标 \(k \neq i\) 成立。

\begin{enumerate}
\def\labelenumi{\arabic{enumi})}
\setcounter{enumi}{10}
\tightlist
\item
  设 A 是赋有豪斯多夫半绝对值的一个环，并设 A 有单位元，且关于由该半绝对值定义的拓扑是完备的。证明 A 的每个极大理想都是闭的（利用命题 13）。
\end{enumerate}

¶ 12) 设 K 是一个非离散的豪斯多夫拓扑除环，\(\varphi: \mathbf{K} \to \mathbf{R}_{+}\) 是一个映射，满足 \(\varphi(0) = 0\) ，\(\varphi(xy) = \varphi(x)\varphi(y)\) ，且使得：若 \(V_{n}\) 表示所有 \(x \in K\) 中满足 \(\varphi(x) \leqslant 1/n\) 者组成的集合，则诸集 \(V_{n}\) 构成 K 中 \(0\) 的一个基本邻域系。

\begin{enumerate}
\def\labelenumi{\alph{enumi})}
\tightlist
\item
  证明存在一个实数 \(a > 0\) ，使得对所有 \(x \in \mathbf{K}\) ，
\end{enumerate}

\[
\varphi (1 + x) \leqslant a (1 + \varphi (x))
\]

{[}否则，将存在 K 的点组成的序列 \((x_{n})\) ，使得 \((1 + x_n)^{-1}\) 与 \(x_{n}(1 + x_{n})^{-1}\) 都趋于零{]}。由此证明：若 \(\psi (x) = (\varphi (x))^{\alpha}\) ，则有

\[
\psi (x + y) \leqslant 2 \sup (\psi (x), \psi (y))
\]

对充分小的 \(\alpha\) 成立。

b) 证明若 \(n = 2^p\) ，则有 \(\psi\left(\sum_{i=1}^{n} x_i\right) \leqslant n \sup (\psi(x_i))\) ，并推出对每个整数 \(m > 0\) 有 \(\psi(m) \leqslant 2m\) 。

\begin{enumerate}
\def\labelenumi{\alph{enumi})}
\setcounter{enumi}{2}
\tightlist
\item
  由 b) 推出：对每个 \(n = 2^p\) 及每个 \(x \in \mathbf{K}\) ，有 \(\psi((1 + x)^{n-1}) \leqslant 2n(1 + \psi(x))^{n-1}\) ，从而 \(\psi\) 是 \(\mathbf{K}\) 上的一个定义 \(\mathbf{K}\) 的拓扑的绝对值。
\end{enumerate}

\begin{enumerate}
\def\labelenumi{\arabic{enumi})}
\setcounter{enumi}{12}
\tightlist
\item
  设 K 是一个非离散的豪斯多夫拓扑除环。设 R 是所有 \(x \in K\) 中满足 \(\lim_{n \to \infty} x^n = 0\) 者组成的集合，N 是集合 \(R \cup R^{-1}\) 的余集。证明：K 上存在一个定义 K 的拓扑的绝对值，当且仅当 (i) R 在 K 中是开的；(ii) 对 K 中 0 的每个邻域 V，存在 K 中 0 的一个邻域 U，使得 \(R.U \subset V\) ；且 (iii) 若 \(x \in R\) 且 \(y \in R \cup N\) ，则 \(yx \in R\) 。
\end{enumerate}

为证明这些条件是充分的，依次证明：a) N 是 K 的非零元素组成的乘法群 \(K^{*}\) 的一个正规子群。

\begin{enumerate}
\def\labelenumi{\alph{enumi})}
\setcounter{enumi}{1}
\item
  在商群 \(K^{*}/N\) 中，规定 \(\dot{x} \leqslant \dot{y}\) ，如果存在 \(x \in \dot{x}\) 与 \(y \in \dot{y}\) 使得 \(yx^{-1} \in R \cup N\) ；证明这个关系是一个与 \(K^{*}/N\) 的群结构相容的序关系，且如此定义的有序群同构于加法群 R 的一个子群（利用第 V 章，§ 3，习题 1）。
\item
  借助习题 12 完成证明。
\end{enumerate}

特别地，非离散局部紧拓扑域的拓扑可以由一个绝对值定义。

\ExerciseSection

\begin{enumerate}
\def\labelenumi{\roman{enumi})}
\tightlist
\item
  \begin{enumerate}
  \def\labelenumii{\alph{enumii})}
  \tightlist
  \item
    构造一个由四个点组成的拓扑空间，它满足公理 \((\mathbf{O}_{\mathbf{V}})\) 但不满足公理 \((\mathbf{O}_{\mathbf{III}})\) 。
  \end{enumerate}
\end{enumerate}

\begin{enumerate}
\def\labelenumi{\alph{enumi})}
\setcounter{enumi}{1}
\item
  满足公理 \((\mathrm{O}_{\mathrm{V}})\) 的拓扑空间 X 也满足公理 \((\mathrm{O}_{\mathrm{III}})\) ，当且仅当它具有下列性质：X 的每个闭子集都是它的诸邻域的交。这时 X 是可一致化的，并且与 X 相伴的完全正则空间（§ 1，习题 4）是正规的。
\item
  若拓扑空间满足公理 (C) 与 \((\mathrm{O}_{\mathrm{III}})\) ，则它满足 \((\mathrm{O}_{\mathrm{V}})\) ，且相伴的完全正则空间是紧的。
\end{enumerate}

\begin{enumerate}
\def\labelenumi{\arabic{enumi})}
\setcounter{enumi}{1}
\tightlist
\item
  设 X 是满足公理 \((\mathbf{O}_{\mathbf{v}})\) 的拓扑空间。
\end{enumerate}

\begin{enumerate}
\def\labelenumi{\alph{enumi})}
\item
  证明 X 的两点 x、y 之间的关系 \(R: \{x\} \cap \{y\} \neq \emptyset\) 等价于关系 \(R' : "f(x) = f(y)\) ——对 X 上每个实值连续函数 f——''；由此证明 R 是 X 上的一个等价关系。
\item
  证明商空间 \(\mathbf{Y} = \mathbf{X} / \mathbf{R}\) 是正规的，且若 \(\varphi: \mathbf{X} \to \mathbf{Y}\) 是典范映射，则每个连续实值函数 \(f\) （在 \(\mathbf{X}\) 上）都可写成 \(f = g \circ \varphi\) 的形式，其中 \(g\) 是 \(\mathbf{Y}\) 上的一个连续实值函数。
\end{enumerate}

\begin{enumerate}
\def\labelenumi{\arabic{enumi})}
\setcounter{enumi}{2}
\tightlist
\item
  若 X 是一个拓扑空间，则下列条件等价：
\end{enumerate}

\begin{enumerate}
\def\labelenumi{\alph{enumi})}
\item
  X 的每个子空间满足公理 \((\mathrm{O}_{\mathrm{V}}^{\prime})\) 。
\item
  X 的每个开子空间满足公理 \((\mathbf{O}_{\mathbf{V}}^{\prime})\) 。
\item
  对 X 的任意两个满足 \(A \cap \overline{B} = B \cap \overline{A} = \varnothing\) 的子集 A 与 B，存在两个不相交的开集 U、V，使得 \(A \subset U\) 且 \(B \subset V\) 。
\end{enumerate}

拓扑空间 \(\mathbf{X}^{-}\) 称为完全正规的，如果它是豪斯多夫的且满足这些条件。每个可度量化空间都是完全正规的。紧空间未必是完全正规的。

\begin{enumerate}
\def\labelenumi{\arabic{enumi})}
\setcounter{enumi}{3}
\tightlist
\item
  每个全序集 X，赋以两个拓扑 \(\mathcal{C}_{+}(\mathbf{X})\) 、\(\mathcal{C}_{-}(\mathbf{X})\) 之一（第 I 章，§ 2，习题 5），都是完全正规的。（若 A 与 B 是 X 的两个满足 \(\mathbf{A} \cap \overline{\mathbf{B}} = \mathbf{B} \cap \overline{\mathbf{A}} = \emptyset\) 的子集，定义邻域 \(V_{x}\) （\(x\) 的，对每个 \(x \in \mathbf{A}\) ），以及邻域 \(W_{y}\) （\(y\) 的，对每个 \(y \in \mathbf{B}\) ），使得 \(V_{x} \cap W_{y} = \emptyset\) 对所有 \(x \in \mathbf{A}\) 及所有 \(y \in \mathbf{B}\) 成立。）
\end{enumerate}

¶ 5) 每个全序集 X，赋以拓扑 \(\mathcal{G}_{0}(\mathbf{X})\) （第 I 章，§ 2，习题 5），都是完全正规的。{[}先考虑 X 是紧的情形；借助第 IV 章，§ 2，习题 6，证明 X 中每个开集都是两两不相交的开区间之并；利用这一结果证明命题，办法是（用习题 3 的记号）考虑闭集 \(\overline{\mathbf{A}}\cap\overline{\mathbf{B}}\) 的余集，然后考虑 \(\overline{\mathbf{B}}\) 的余集；为过渡到 X 任意的普遍情形，利用第 IV 章，§ 4，习题 7。{]}

\begin{enumerate}
\def\labelenumi{\arabic{enumi})}
\setcounter{enumi}{5}
\item
  证明：在正规空间 X 中，X 的每个是闭集的可数并的子空间 Y 都是正规的。{[}设 A、B 是 Y 的两个不相交的闭子集，并设 Y 是闭子集组成的一个递增序列 \((\mathbf{Y}_n)\) （在 \(\mathbf{X}\) 中）的并。用归纳法定义开子集的两个序列 \((\mathbf{U}_n), (\mathbf{V}_n)\) （\(\mathbf{X}\) 的），使得：(i) \(\overline{\mathbf{U}_n \cap \mathbf{Y}_n}\) 与 \(\overline{\mathbf{V}_n \cap \mathbf{Y}_n}\) 不相交；(ii) \(\mathbf{U}_n \cap \mathbf{Y}_n\) 包含 \(\mathbf{A} \cap \mathbf{Y}_n\) 以及诸集 \(\overline{\mathbf{U}_i \cap \mathbf{Y}_i}\) （对 \(1 \leqslant i \leqslant n\) ）；且 \(\mathbf{V}_n \cap \mathbf{Y}_n\) 包含 \(\mathbf{B} \cap \mathbf{Y}_n\) 以及诸 \(\overline{\mathbf{V}_i \cap \mathbf{Y}_i}, 1 \leqslant i \leqslant n\) 。{]}
\item
  拓扑空间 X 称为完备正规的，如果它是正规的，且 X 的每个闭子集都是可数个开集的交（或者等价地，X 的每个开子集都是可数个闭集的并）。
\end{enumerate}

\begin{enumerate}
\def\labelenumi{\alph{enumi})}
\item
  豪斯多夫空间 X 是完备正规的，当且仅当对 X 的任一闭子集 A，存在 X 上的一个实值连续函数 f，使得 \(\overline{f}^{1}(0)=A\) {[}利用 § 1，习题 16 a) 的方法{]}。b) 证明每个完备正规空间 X 都是完全正规的，且 X 的每个子空间都是完备正规的{[}利用习题 3 b) 与 6{]}。
\item
  完备正规空间上的每个下半连续实值函数都是连续函数组成的一个序列的上包络（利用 § 2，第 7 号，命题 11 的方法）。
\item
  在不可数离散空间上加一个无穷远点所得的紧空间 X 是完全正规但非完备正规的，且 X 的每个子空间都是仿紧的。
\end{enumerate}

§ 8) 证明 § 2，习题 13 a) 中定义的非可度量化紧空间 X 是完备正规的，而 § 2，习题 13 d) 中定义的非可度量化紧空间 Y 是完全正规的但不是完备正规的。

¶ 9) a) 设 X 是一个仿紧空间，其中每个点都有可数的基本邻域系，Y 是一个可数紧的正规空间（§ 2，习题 14）。证明 X × Y 是正规的。{[}设 A、B 是 X × Y 的两个不相交的闭子集；对每个 x ∈ X，证明存在 x 在 X 中的一个邻域 U\_x 及 Y 中两个开集 V\_x、W\_x，使得 \(\overline{V}_x \cap \overline{W}_x = \emptyset\) ，且对每个 z ∈ U\_x 有 A(z) ⊂ V\_x 及 B(z) ⊂ W\_x；为证这一点，用反证法论证。设 (T\_λ)\_\{λ ∈ L\} 是 X 的一个局部有限开覆盖，它加细覆盖 (U\_x)\_\{x ∈ X\}，并设 (f\_λ) 是从属于 (T\_λ) 的一个单位分解；对每个 λ ∈ L，取 x\_λ 使 T\_λ ⊂ U\_\{x\_λ\}，并令 g\_λ 为 Y 到 {[}0, 1{]} 中的一个连续映射，它在 \(\overline{V}_{x_λ}\) 上等于 0，在 \(\overline{W}_{x_λ}\) 上等于 1；考虑 X × Y 上的函数 \(\sum_{\lambda} f_\lambda(x) g_\lambda(y)\) 。{]}

\begin{enumerate}
\def\labelenumi{\alph{enumi})}
\setcounter{enumi}{1}
\tightlist
\item
  设 Y = {[}a, b{[} 是第 I 章，§ 9，习题 12 中定义的局部紧空间，X = {[}a, b{]} 是在 Y 上加一个无穷远点所得的紧空间。Y 是正规的（习题 4）且可数紧（§ 2，习题 15），但乘积 X×Y 不是正规的（利用第 II 章，§ 4，习题 4）。
\end{enumerate}

¶ 10) 设 L 是一个不可数集，X 表示完全正则空间 \(\mathbf{N}^{\mathrm{L}}\) （N 赋以离散拓扑）。设 A（相应地 B）是 X 中所有满足如下条件的 \(x = (x(\lambda))_{\lambda \in \mathrm{L}}\) 组成的集合：对每个整数 \(k \neq 0\) （相应地 \(k \neq 1\) ），由 \(\lambda \in \mathbf{L}\) 中满足 \(x(\lambda) = k\) 者组成的集至多有一个元素。a) 证明 A 与 B 是 X 的两个不相交的闭子集。

\begin{enumerate}
\def\labelenumi{\alph{enumi})}
\setcounter{enumi}{1}
\item
  设 U、V 是 X 的两个开子集，满足 \(A \subset U\) 与 \(B \subset V\) 。令 \(\Phi\) 表示 X 的所有初等集（第 I 章，§ 4，第 1 号）中其异于 N 的投影均由单独一个元素组成者所成的集合；对每个 \(W \in \Phi\) ，令 \(\mathrm{H}(\mathrm{W})\) 为所有满足下列条件的 \(\lambda \in L\) 组成的集合：\(\mathrm{pr}_{\lambda}(\mathrm{W})\) 由单独一个元素组成。证明存在 A 的点组成的一个序列 \((x_{n})\) 、一个集序列 \((\mathbf{U}_{n})\) （集取自 \(\Phi\) ）、L 的不同元素组成的一个序列 \((\lambda_{n})\) 以及整数的一个严格递增序列 \((m(n))_{n \in \mathbb{N}}\) ，它们具有下列性质：(i) \(U_{n} \subset U\) 是 \(x_{n}\) 的一个邻域；(ii) \(\mathrm{H}(\mathrm{U}_{n})\) 是诸 \(\lambda_{k}\) （\(k \leqslant m(n)\) ）组成的集合；(iii) \(x_{0}(\lambda) = 0\) 对每个 \(\lambda \in L\) 成立，且对每个 n \textgreater{} 0，\(x_{n}(\lambda_{k}) = k\) （若 \(k \leqslant m(n - 1)\) ），而 \(x_{n}(\lambda) = 0\) 对所有其他 \(\lambda \in L\) 成立。
\item
  设 \(y \in \mathbf{B}\) 是这样的点：\(y(\lambda_k) = k\) 对每个整数 \(k\) 成立，且 \(y(\lambda) = \mathbf{1}\) 对所有 \(\lambda \in \mathbf{L}\) 中异于诸 \(\lambda_k\) 者成立。设 \(\mathbf{V}_0 \subset \mathbf{V}\) 是 \(\Phi\) 的一个包含 \(y\) 的集合，并设 \(n\) 是使 \(\lambda_k \in \mathbf{L} - \mathbf{H}(\mathbf{V}_0)\) （对所有 \(k > m(n)\) ）的整数。证明 \(\mathbf{U}_{n+1} \cap \mathbf{V}_0 \neq \emptyset\) ，并推出 \(\mathbf{X}\) 不是正规的。
\end{enumerate}

\begin{enumerate}
\def\labelenumi{\arabic{enumi})}
\setcounter{enumi}{10}
\tightlist
\item
  由习题 10 推出：
\end{enumerate}

\begin{enumerate}
\def\labelenumi{\alph{enumi})}
\item
  若非空豪斯多夫空间的一个乘积 \(\prod_{i\in I}X_i\) 是正规的，则除去一个可数指标集外，\(X_i\) 都是可数紧的（§ 2，习题 14）。
\item
  非空可度量化空间的乘积 \(\prod_{i\in I}X_i\) 是正规的，当且仅当除去一个可数指标集外 \(X_{i}\) 都是紧的；此时乘积空间 \(\prod_{i\in I}X_{i}\) 是仿紧的。
\end{enumerate}

\begin{enumerate}
\def\labelenumi{\arabic{enumi})}
\setcounter{enumi}{11}
\tightlist
\item
  设 X、Y 是两个豪斯多夫空间。
\end{enumerate}

\begin{enumerate}
\def\labelenumi{\alph{enumi})}
\item
  假设 X 中有一个闭集 A，它不是可数个开集的交，且 Y 有一个可数无限子集 B，它不是闭的。取 \(b \in \overline{B} - B\) ，并考虑子集 \(C = A \times B\) 与 \(D = (X - A) \times \{b\}\) （在 \(X \times Y\) 中）。证明 \(C \cap \overline{D} = \overline{C} \cap D = \varnothing\) ，但不存在 \(X \times Y\) 的开子集的对 U、V，使得 \(C \subset U\) 、\(D \subset V\) 且 \(U \cap V = \varnothing\) 。
\item
  若 \(X \times Y\) 是完全正规的，则或者 (i) 空间 X、Y 之一是完备正规的，或者 (ii) 空间 X、Y 之一具有每个可数无限子集都闭的性质（参看 § 5，习题 16）。
\item
  设 \(X = [a, b]\) 是一个不可数的良序集，使得 \([a, x]\) 对每个 x \textless{} b 都是可数的。给 X 赋以如下拓扑：\(\{x\}\) 对每个 x \textless{} b 是开的，且诸区间 \([x, b]\) （x \textless{} b）构成 b 的一个基本邻域系。证明 \(X \times X\) 是完全正规的，但 X 不是完备正规的。
\item
  证明：若 \(\mathbf{X}\) 是紧空间，且 \(\mathbf{X} \times \mathbf{X} \times \mathbf{X}\) 完全正规，则 \(\mathbf{X}\) 可度量化（利用 § 2，第 4 号，定理 1）。
\end{enumerate}

¶ 13) 设 \((\mathbf{X}_{t})_{t\in I}\) 是豪斯多夫拓扑空间组成的一个不可数族，其中每个空间都不止含一个点。对每个 \(t\in I\) ，设 \(a_{t}\) 与 \(b_{t}\) 是 \(\mathbf{X}_{t}\) 的两个不同的点。设 \(\mathbf{X}\) 是积空间 \(\prod_{t\in I}\mathbf{X}_{t}\) ，并设 Y 是 X 中由这样的点 \((x_{t})_{t\in I}\) 组成的子空间：除去一个可数的指标集外，均有 \(x_{t}=a_{t}\) ；又设 b 是 X 的点 \(b_{t}\) 。

在积空间 \(\mathbf{X} \times \mathbf{Y}\) 中，设 \(\mathbf{A}\) 是 \(\mathbf{Y} \times \mathbf{Y}\) 的对角线，并设 \(\mathbf{B}\) 是集 \(\{b\} \times \mathbf{Y}\) 。证明 \(\mathbf{A}\) 与 \(\mathbf{B}\) 是闭集，并且不存在开集对 \(\mathbf{U}, \mathbf{V}\) （在 \(\mathbf{X} \times \mathbf{Y}\) 中），满足 \(\mathbf{A} \subset \mathbf{U}, \mathbf{B} \subset \mathbf{V}\) 以及 \(\mathbf{U} \cap \mathbf{V} = \emptyset\) 。{[}设 \(\mathbf{U}, \mathbf{V}\) 是两个开集，满足 \(\mathbf{A} \subset \mathbf{U}\) 与 \(\mathbf{B} \subset \mathbf{V}\) ；证明存在一个递增序列 \((\mathbf{H}_n)\) （由 \(\mathbf{I}\) 的可数子集组成），使得：若 \(x_n\) 表示 \(\mathbf{Y}\) 中这样的点——其指标 \(i \in \mathbf{H}_n\) 的每个坐标等于 \(b_i\)，而其指标 \(i \notin \mathbf{H}_n\) 的每个坐标等于 \(a_i\) ——则点 \((x_{n+1}, x_n)\) 属于 \(\mathbf{V}\) （对每个整数 \(n\) ）；证明序列 \((x_{n+1}, x_n)\) 趋于 \(\mathbf{A}\) 的一个点，并由此推出 \(\mathbf{U} \cap \mathbf{V} \neq \emptyset\) 。{]}

由此构造一个非正规的连通拓扑环的例子。

并由此推出：由豪斯多夫拓扑空间组成的一个不可数族（其中每个空间至少有两个不同的点）的积不可能是完全正规的（在诸 \(X_{t}\) 都是两个点的离散空间时利用上述结果）。

\begin{enumerate}
\def\labelenumi{\arabic{enumi})}
\setcounter{enumi}{13}
\item
  设 X 是一个非正规的豪斯多夫空间，并设 A, B 是 X 的两个不相交的闭子集，使得不存在 X 中满足关系 \(A \subset U\) 与 \(B \subset V\) 的不相交开集对 U, V。设 R 是 X 上的等价关系，其等价类为：集 A、集 B、以及诸集 \(\{x\}\) （其中 \(x \in \complement(A \cup B)\) ）。证明 R 的图象在 \(X \times X\) 中是闭的，且关系 R 是闭的，但商空间 X/R 不是豪斯多夫空间（参看第 I 章，§ 8，第 3 号，命题 8）。
\item
  \begin{enumerate}
  \def\labelenumii{\alph{enumii})}
  \tightlist
  \item
    设 X 是正规空间，R 是 X 上的一个闭等价关系。证明商空间 X/R 是正规的（利用第 I 章，§ 5，第 4 号，命题 10）。
  \end{enumerate}
\end{enumerate}

\begin{enumerate}
\def\labelenumi{\alph{enumi})}
\setcounter{enumi}{1}
\item
  设 X 是一个局部紧的 σ 紧空间，R 是 X 上的等价关系，其图象在 X×X 中是闭的。证明 X/R 是正规的（参看第 I 章，§ 10，习题 19）。
\item
  设 X 是 R 中由形如 1/n（其中 n 是异于 0 与 ±1 的整数）的点组成的集的余集。在可度量化空间 X 上考虑这样的等价关系 R：每个非整数的 \(x \in X\) 的等价类由点 x 与 1/x 组成，而每个整数的等价类由该整数单独组成。证明 R 是开的，且 R 的图象在 \(X \times X\) 中是闭的，但 X/R 不是正则的。
\end{enumerate}

\begin{enumerate}
\def\labelenumi{\arabic{enumi})}
\setcounter{enumi}{15}
\tightlist
\item
  对每个覆盖 \(\Re\) （拓扑空间 \(X\) 的），用 \(V_{\Re}\) 表示诸集 \(U \times U\) （\(U \in \Re\)）的并。覆盖 \(\Re\) （\(X\) 上的）称为均匀的，如果有一个邻域 \(W\) （对角线 \(\Delta\) 的，\(X \times X\) 中的），使得覆盖 \((W(x))_{x \in X}\) 加细 \(\Re\)。称 \(\Re\) 为可除的，如果有一个邻域 \(W\) （\(\Delta\) 的，在 \(X \times X\) 中），使得 \(\dot{W} \subset V_{\Re}\) 。 a) 证明 \(X\) 的每个均匀开覆盖都是可除的{[}参看习题 19 b){]}。
\end{enumerate}

\begin{enumerate}
\def\labelenumi{\alph{enumi})}
\setcounter{enumi}{1}
\tightlist
\item
  设 \(\Re\) 是 \(X\) 的一个开覆盖，使得存在一个局部有限闭覆盖 \(\mathfrak{S}\) （\(X\) 的）加细 \(\Re\)。证明 \(\Re\) 是均匀的。（对每个集 \(A \in \mathfrak{S}\)，设 \(U_A\) 是 \(\Re\) 中包含 \(A\) 的一个集，并设 \(V_A\) 是 \(X \times X\) 中 \(U_A \times U_A\) 与 \((X - A) \times (X - A)\) 的并；考虑集 \(W = \bigcap_{A \in \mathfrak{S}} V_A\) 。）
\end{enumerate}

¶ 17) a) 若拓扑空间 X 的每个有限开覆盖都可除（习题 16），则 X 满足公理 \((\mathrm{O}_{\mathbf{V}}^{\prime})\) （考虑 X 的一个二元开覆盖，即由 X 的两个子集所组成的覆盖）。反之，若 X 满足 \((\mathrm{O}_{\mathbf{V}}^{\prime})\) ，则 X 的每个有限开覆盖都是均匀的{[}利用第 3 号的定理 3（只要 X 满足 \((\mathrm{O}_{\mathbf{V}}^{\prime})\) ，它便成立），先化归到二元覆盖的情形，然后再处理这一情形{]}。当 R 跑遍 X 的有限开覆盖所成的集时，诸集 \(V_{R}\) 构成 X 上一个一致结构的围域的基本系，该一致结构与 X 的拓扑相容；这个一致结构称为``有限开覆盖的一致结构''。

\begin{enumerate}
\def\labelenumi{\alph{enumi})}
\setcounter{enumi}{1}
\item
  假设 X 是正规的。证明有限开覆盖的一致结构 U 与如下的最粗一致结构 \(U'\) 相同：X 到 {[}0, 1{]} 的所有连续映射关于它都一致连续{[}即 X 上由其斯通–切赫紧化（§ 1，习题 7）的一致结构所诱导的一致结构{]}。{[}为证 U 比 \(U'\) 粗，利用第 3 号的定理 3 与公理 \((\mathrm{O_V})\)，在 \(\mathbf{X}\) 上构造形如 \((x, y) \to |f(x) - f(y)|\) 的伪度量组成的一个有限族，使得由这些伪度量所定义的一致结构的一个围域含于一致结构 \(\mathfrak{U}]\) 的一个围域。
\item
  设 X 是正规空间，并设 \(\beta\mathbf{X}\) 是它的斯通–切赫紧化。对 X 的每个闭子集 Y，设 \(\overline{\mathbf{Y}}\) 是 Y 在 \(\beta\mathbf{X}\) 中的闭包。证明由 Y 的恒等映射扩张而得的、从 \(\beta\mathbf{Y}\) 到 \(\overline{\mathbf{Y}}\) 的连续映射（\(\beta\mathbf{Y}\) 是 Y 的斯通–切赫紧化）是一个同胚{[}或者利用 \(b\) )，或者利用定理 2 与 \(\beta\mathbf{Y}]\) 的万有性质。若 Y 与 Z 是 X 的两个闭子集，证明 \(\overline{\mathbf{Y}} \cap \overline{\mathbf{Z}} = \overline{\mathbf{Y} \cap \mathbf{Z}}\) 在 \(\beta\mathbf{X}\) 中成立（若 \(x_0 \notin \overline{\mathbf{Y}} \cap \overline{\mathbf{Z}}\) 且 \(x_0 \in \overline{\mathbf{Y}}\) ，考虑 \(x_0\) 在 \(\beta\mathbf{X}\) 中的一个闭邻域 V，使得 V \(\cap\) Y 与 Z 是 X 的不相交闭子集）。
\end{enumerate}

¶ 18) 拓扑空间 X 的子集组成的一个族 \((A_{\alpha})\) 称为离散的，如果对每个 \(x \in X\) ，存在 x 的一个邻域，它至多与族中的一个集 \(A_{\alpha}\) 相交。豪斯多夫空间 X 称为集体正规的，如果对 X 的闭子集组成的每个离散族 \((A_{\alpha})\) ，存在 X 的两两不相交的开子集组成的一个族 \((U_{\alpha})\) ，使得 \(A_{\alpha} \subset U_{\alpha}\) 对每个指标 \(\alpha\) 成立。每个集体正规空间都是正规的。

\begin{enumerate}
\def\labelenumi{\alph{enumi})}
\item
  完全正则空间 X 的每个开覆盖都可除（习题 16），当且仅当对角线 \(\Delta\) （在 \(\mathbf{X} \times \mathbf{X}\) 中）的邻域所成之集是 X 的万有一致结构的围域滤子（ \(\S\) 1，习题 5）。证明：若此条件满足，则 X 是集体正规的{[}若 \((\mathbf{A}_{\alpha})\) 是 X 的闭子集组成的一个离散族，考虑开覆盖 \((\mathbf{V}_{\alpha})\) ，其中 \(\mathbf{V}_{\alpha} = \mathbf{X} - \bigcup_{\beta \neq \alpha} \mathbf{A}_{\beta}\) 对每个 \(\alpha\) 成立{]}。
\item
  设 Y = {[}a, b{]} 是第 I 章，§ 9，习题 12 中定义的局部紧空间；设 Y₀ 是赋以离散拓扑的集 Y，并设 Z 是集 Y₀ ∪ \{b\}，其中 Y₀ 的点是开集，而诸集 {]}x, b{[}（x 跑遍 Y₀）构成 b 的一个基本邻域系。证明空间 X = Y × Z 是集体正规的（回忆 Y 的子集组成的任何无限族都不可能是离散的，并利用习题 4）。设 R 是 X 的开覆盖，由 Y × Y₀ 以及诸积 {[}a, x{]} × {]}x, b{]} （其中 x \textless{} b）组成；证明 R 不是可除的，并且对角线 Δ （在 X × X 中的）的邻域所成之集因而不是 X 上一个一致结构的围域滤子{[}利用第 I 章，§ 9，习题 12 a){]}。
\end{enumerate}

¶ 19) a) 正则空间 X 是仿紧的，当且仅当 X 的每个开覆盖都是均匀的（习题 16）。{[}为证必要性，利用引理 5 与习题 16 b)；为证充分性，利用 § 4 的习题 16 a)、§ 1 第 4 号的命题 2 以及 § 4 第 5 号的定理 4{]}。

\begin{enumerate}
\def\labelenumi{\alph{enumi})}
\setcounter{enumi}{1}
\item
  给出集体正规但非仿紧的空间上的一个可除而非均匀的覆盖的例子{[}利用 a) 与第 II 章，§ 4，习题 4{]}。
\item
  证明仿紧空间 X 关于其万有一致结构（§ 1，习题 5）是完备的。{[}若关于这个一致结构的一个柯西滤子 F 没有触点，则每个点 \(x \in X\) 都有一个开邻域 \(V_{x}\) ，它至少不与 F 的一个集相交；利用 a) 与习题 18 a)。{]}
\end{enumerate}

\begin{enumerate}
\def\labelenumi{\arabic{enumi})}
\setcounter{enumi}{19}
\tightlist
\item
  \begin{enumerate}
  \def\labelenumii{\alph{enumii})}
  \tightlist
  \item
    正则空间 X 是仿紧的，如果对 X 的每个开覆盖 R，存在一个序列 \((\mathfrak{S}_{n})\) ，使得 \(\mathfrak{S} = \bigcup_{n} \mathfrak{S}_{n}\) 是 \(\mathbf{X}\) 的一个加细 \(R\) 的开覆盖，且每个 \(\mathfrak{S}_n\) 都是局部有限族（参看引理 4、5、6）。
  \end{enumerate}
\end{enumerate}

\begin{enumerate}
\def\labelenumi{\alph{enumi})}
\setcounter{enumi}{1}
\item
  由 a) 推出：在仿紧空间中，闭集组成的每个可数并都是一个仿紧子空间（参看 § 5，习题 15）。
\item
  在正则空间 X 中，设 F 是闭集组成的一个局部有限族，其中每个闭集都是 X 的仿紧子空间。证明 F 中诸集的并是 X 的一个仿紧子空间（利用第 5 号的引理 6；参看 § 5，习题 15）。
\item
  证明：若 X 是仿紧空间，而 Y 是紧集组成的一个可数并的正则空间，则 \(X \times Y\) 是仿紧的{[}把 Y 嵌入它的斯通–切赫紧化并利用 b)；参看 § 5，习题 16{]}。
\end{enumerate}

¶ 21) a) 设 X 是仿紧空间，R 是 X 上的一个闭等价关系，使得关于 R 的每个等价类都是紧的。证明 X/R 是仿紧的{[}利用习题 15 a)、第 3 号的定理 3 以及第 5 号的引理 6；参看习题 15 c){]}。

\begin{enumerate}
\def\labelenumi{\alph{enumi})}
\setcounter{enumi}{1}
\item
  设 X 是豪斯多夫空间，R 是 X 上的一个闭等价关系，使得关于 R 的每个等价类都是紧的。证明：若 X/R 仿紧，则 X 也仿紧（利用第 I 章，§ 9，第 10 号，命题 17 的论证）。
\item
  设 Y 是局部紧但非仿紧的空间（参看第 I 章，§ 9，习题 12），并设 \(Y_{0}\) 是 Y 的一点紧化。设 X 是 Y 与 \(Y_{0}\) 的拓扑和，并设 R 是 X 上的等价关系，它把 Y 的每个点与其在 \(Y_{0}\) 中的典范象等同起来。关系 R 是开的，关于 R 的每个等价类至多含两个点，且 X/R 同胚于 \(Y_{0}\) ，因而是紧的。
\end{enumerate}

\begin{enumerate}
\def\labelenumi{\arabic{enumi})}
\setcounter{enumi}{21}
\item
  正则空间 \(\mathbf{X}\) 可度量化，当且仅当存在一个序列 \((\mathfrak{B}_n)\) （由 \(\mathbf{X}\) 的开子集的局部有限族组成），使得 \(\mathfrak{B} = \bigcup_{n} \mathfrak{B}_n\) 是 \(\mathbf{X}\) 的拓扑的一个基。{[}为证条件是必要的，利用第 5 号的引理 3。为证它是充分的，先用引理 4、5、6 证明 X 是仿紧的。对每对整数 m, n 以及每个 \(U \in \mathfrak{B}_{m}\) ，设 \(U'\) 是诸集 \(V \in \mathfrak{B}_{n}\) （其中 \(\overline{V} \subset U\) ）的并。证明 \(\overline{U}' \subset U\) 。设 \(f_{U}: X \to [0, 1]\) 是一个连续映射，它在 X---U 上等于 0，在 \(U'\) 上等于 1；并设 \(d_{mn}(x, y) = \sum_{U \in \mathfrak{B}_{m}} |f_{U}(x) - f_{U}(y)|\) ；证明诸伪度量 \(d_{mn}\) 定义 X 的拓扑。{]}（``长田–斯米尔诺夫定理''）。特别地，具有可数基的每个正则空间都是可度量化的。
\item
  \begin{enumerate}
  \def\labelenumii{\alph{enumii})}
  \tightlist
  \item
    每个正则的林德勒夫空间（第 I 章，§ 9，习题 15）都是仿紧的{[}利用习题 20 a){]}。
  \end{enumerate}
\end{enumerate}

\begin{enumerate}
\def\labelenumi{\alph{enumi})}
\setcounter{enumi}{1}
\tightlist
\item
  林德勒夫空间与紧空间的积是林德勒夫空间（按第 I 章，§ 9，第 10 号，命题 17 的证明方式论证；参看 § 5，习题 16）。
\end{enumerate}

\begin{enumerate}
\def\labelenumi{\arabic{enumi})}
\setcounter{enumi}{23}
\tightlist
\item
  \begin{enumerate}
  \def\labelenumii{\alph{enumii})}
  \tightlist
  \item
    设 X 是可度量化空间，R 是 X 上的一个闭等价关系，其所有等价类都是紧的。证明 X/R 是可度量化的。{[}对 X 的由关于 R 饱和的集组成的开覆盖证明第 5 号引理 3 的类似命题，以此来应用长田–斯米尔诺夫定理（习题 22）；用 \(F_{n\alpha}\) 表示含于由所有 \(x \in X\) （满足 \(d(x, X - U_{\alpha}) > 2^{-n}\) ）组成之集的最大饱和开集；用 \(F_{n\alpha}'\) 表示由所有 \(x \in X\) （满足 \(d(x, X - U_{\alpha}) \geqslant 2^{-n}\) ）组成之集的饱和化；并用 \(G_{n\alpha}\) 表示由所有 \(x \in F_{n\alpha}\) （满足 \(x \notin F_{n+1,\beta}'\) ，对每个 \(\beta < \alpha\) ）组成的集。{]}
  \end{enumerate}
\end{enumerate}

\begin{enumerate}
\def\labelenumi{\alph{enumi})}
\setcounter{enumi}{1}
\item
  把 a) 的结果推广到如下情形：关系 R 是闭的，且 X/R 的每个点都有可数的基本邻域系{[}利用 § 2 的习题 20 b){]}。
\item
  设 X 是拓扑空间，\((\mathbf{F}_{\alpha})\) 是 X 的一个局部有限闭覆盖。证明：若诸子空间 \(F_{\alpha}\) 的每一个都可度量化，则 X 可度量化（把 X 看作诸 \(F_{\alpha}\) 的拓扑和的一个商空间）。
\item
  仿紧空间（其中每个点都有一个可度量化的邻域）是可度量化的{[}应用 c)；参看第 I 章，§ 9，习题 12 b)，以及第 IX 章，§ 4，习题 25 d)、§ 2，习题 15 与 § 5，习题 15{]}。
\end{enumerate}

\begin{enumerate}
\def\labelenumi{\arabic{enumi})}
\setcounter{enumi}{24}
\tightlist
\item
  豪斯多夫空间 X 称为亚紧的，如果对 X 的任何开覆盖 R，存在 X 的一个加细 R 的逐点有限开覆盖。
\end{enumerate}

\begin{enumerate}
\def\labelenumi{\alph{enumi})}
\item
  证明亚紧空间的每个闭子空间都是亚紧的，并且若亚紧空间 X 的每个开子空间都亚紧，则 X 的每个子空间都亚紧。
\item
  设 X 是豪斯多夫空间，R 是 X 上的一个闭等价关系，其所有等价类都是紧的。证明：若 X/R 亚紧，则 X 亚紧{[}按习题 21 b) 的方式论证{]}。
\item
  证明既亚紧又可数紧（§ 2，习题 14）的拓扑空间是紧的。{[}利用佐恩引理证明每个逐点有限开覆盖都包含一个极小开覆盖，并利用 § 2，习题 14 a)。{]}
\end{enumerate}

\(d)\) 第 I 章，§ 9，习题 12 中定义的局部紧完全正规空间 \(\mathbf{X} = [a, b[\) 不是亚紧的（*）。

\begin{enumerate}
\def\labelenumi{\alph{enumi})}
\setcounter{enumi}{4}
\tightlist
\item
  设 Y 是可度量化空间，A 是 Y 的一个子集，使得 A 与其余集都在 Y 中稠密。用 X 表示给 Y 赋以由 Y 的开集与集 A 生成的拓扑所得的非正则空间。证明 X 是亚紧的。（对每个 \(x \in X\) ，考虑 x 的一个邻域 \(V_{x}\) ，它含于 X 的给定开覆盖 R 的某个集中，使得 \(V_{x} \subset A\) （若 \(x \in A\) ），且使得 \(V_{x}\) 是 x 在 Y 中的一个邻域（若 \(x \notin A\) ）。注意 X 的子空间 A 是仿紧的，而且诸 \(V_{x}\) （\(x \notin A\) ）的并也是仿紧的。）
\end{enumerate}

\begin{enumerate}
\def\labelenumi{\arabic{enumi})}
\setcounter{enumi}{25}
\tightlist
\item
  \begin{enumerate}
  \def\labelenumii{\alph{enumii})}
  \tightlist
  \item
    证明每个伪紧的正规空间 X（§ 1，习题 22）都是可数紧的（§ 2，习题 14）。{[}证明：若 \((x_n)\) 是 X 的互异点组成的、没有触点的一个序列，则 X 上存在一个连续的有限实值函数 f，使得 \(f(x_n) = n\) 对所有 n 成立。{]}
  \end{enumerate}
\end{enumerate}

\begin{enumerate}
\def\labelenumi{\alph{enumi})}
\setcounter{enumi}{1}
\tightlist
\item
  设 \(Y = [a, b[\) 是第 I 章，§ 9，习题 12 中定义的局部紧空间，并设 \(Y_{0}\) 是由 Y 添加一个无穷远点（它可与 b 等同）而得的紧空间。设 c 是 Y 中使 {[}a, c{[} 为无限集的最小元素，并设 Z 是 Y 的紧子空间 {[}a, c{]} 。设 X 是紧积空间 \(Y_{0} \times Z\) 中除去点 (b, c) 后的余集。证明 X 局部紧但非正规{[}参看习题 12 a){]}；X 伪紧但非可数紧。给出 X 的一个非伪紧的闭子空间的例子。再给出 X 上的一个不达到其下确界的下半连续函数的例子{[}参看 § 2，习题 14 i){]}。
\end{enumerate}

\begin{enumerate}
\def\labelenumi{\arabic{enumi})}
\setcounter{enumi}{26}
\tightlist
\item
  拓扑空间 X 称为局部仿紧的，如果 X 的每个点都有一个闭邻域，它是 X 的一个仿紧子空间。
\end{enumerate}

\begin{enumerate}
\def\labelenumi{\alph{enumi})}
\item
  证明局部仿紧空间是完全正则的。
\item
  在习题 26 b) 中定义的紧空间 \(Y_{0} \times Z\) 中，设 H 是由点 \((b, y)\) （其中 y \textless{} c）组成的集的余集。证明 H 正规但非局部仿紧{[}参看第 I 章，§ 9，习题 12 b){]}。
\item
  设 X 是拓扑空间，\(\Phi\) 是 X 的一个由 X 的仿紧子空间闭集组成的覆盖，使得每个 \(A \in \Phi\) 都在 X 中有一个属于 \(\Phi\) 的邻域。设 \(X'\) 是 X 与一个点 \(\omega\) 组成的和集；在 \(X'\) 上如下定义拓扑：取 x 在 X 中的全体邻域所成之集，作为 \(x \in X\) 在 \(X'\) 中的基本邻域系，而 \(\omega\) 在 \(X'\) 中的基本邻域系则取为由 \(\Phi\) 的集的余集所生成的滤子基{[}参看习题 20 c){]}。证明 \(X'\) 是仿紧的。
\end{enumerate}

\begin{enumerate}
\def\labelenumi{\arabic{enumi})}
\setcounter{enumi}{27}
\tightlist
\item
  设 X 是度量空间，d 是 X 上的度量，A 是 X 的一个闭子集，并设 \(f: A \to [1, 2]\) 是连续实值函数。证明 X 上的实值函数 g——\(g(x) = f(x)\) （若 \(x \in A\) ），而
\end{enumerate}

\[
g (x) = \frac {\mathrm{1}}{d (x , \mathrm{A})} \inf _ {y \in \mathrm{A}} (f (y) d (x, y))
\]

若 \(x \in \mathbb{C}A\) 则取上式——在 \(X\) 上连续。由此对度量空间给出定理 2 的一个证明。

\begin{enumerate}
\def\labelenumi{\arabic{enumi})}
\setcounter{enumi}{28}
\item
  设 X 是豪斯多夫拓扑空间。则 X 正规，当且仅当：对 X 的每个闭子集 A 以及每个在 A 上定义且关于诱导拓扑连续的实值函数 f，存在 X 上的一个连续伪度量 d，f 关于它一致连续。（为证条件是充分的，证明每个定义在闭子集 A 上且关于诱导拓扑连续的实值函数 f 都能扩张为 X 上的连续函数；为此，考虑与由伪度量 d 在 X 上定义的一致结构相伴的豪斯多夫空间 \(\overline{X}\) 。）
\item
  设 X 是正规空间，g（相应地，f）是 X 上的上（相应地，下）半连续函数，满足 \(g \leqslant f\) 。证明 X 上存在连续实值函数 h，使得 \(g \leqslant h \leqslant f\) 。{[}化归到 f 与 g 取值于 {[}0, 1{]} 的情形。对每个二进数 \(r \in [0, 1]\) ，设 \(\mathrm{F}(r)\) {[}相应地，\(\mathrm{G}(r)\) {]} 为由所有 \(x \in X\) （满足 \(f(x) \leqslant r\) ，相应地，\(g(x) < r\) ）组成的集。用归纳法定义开集的一个族 \(\mathrm{U}(r)\) （r 跑遍含于 {[}0, 1{]} 的二进数所成的集），使得当 \(r < r'\) 时有 \(\mathrm{F}(r) \subset \mathrm{U}(r')\) 、\(\overline{\mathrm{U}(r)} \subset \mathrm{U}(r')\) 与 \(\overline{\mathrm{U}(r)} \subset \mathrm{G}(r')\) ，办法是模仿第 1 号的定理 1 的证明；然后像定理 1 那样完成证明。{]}
\end{enumerate}
